% Options for packages loaded elsewhere
\PassOptionsToPackage{unicode}{hyperref}
\PassOptionsToPackage{hyphens}{url}
\PassOptionsToPackage{dvipsnames,svgnames,x11names}{xcolor}
%
\documentclass[
]{scrartcl}

\usepackage{amsmath,amssymb}
\usepackage{iftex}
\ifPDFTeX
  \usepackage[T1]{fontenc}
  \usepackage[utf8]{inputenc}
  \usepackage{textcomp} % provide euro and other symbols
\else % if luatex or xetex
  \usepackage{unicode-math}
  \defaultfontfeatures{Scale=MatchLowercase}
  \defaultfontfeatures[\rmfamily]{Ligatures=TeX,Scale=1}
\fi
\usepackage{lmodern}
\ifPDFTeX\else  
    % xetex/luatex font selection
\fi
% Use upquote if available, for straight quotes in verbatim environments
\IfFileExists{upquote.sty}{\usepackage{upquote}}{}
\IfFileExists{microtype.sty}{% use microtype if available
  \usepackage[]{microtype}
  \UseMicrotypeSet[protrusion]{basicmath} % disable protrusion for tt fonts
}{}
\makeatletter
\@ifundefined{KOMAClassName}{% if non-KOMA class
  \IfFileExists{parskip.sty}{%
    \usepackage{parskip}
  }{% else
    \setlength{\parindent}{0pt}
    \setlength{\parskip}{6pt plus 2pt minus 1pt}}
}{% if KOMA class
  \KOMAoptions{parskip=half}}
\makeatother
\usepackage{xcolor}
\setlength{\emergencystretch}{3em} % prevent overfull lines
\setcounter{secnumdepth}{5}
% Make \paragraph and \subparagraph free-standing
\ifx\paragraph\undefined\else
  \let\oldparagraph\paragraph
  \renewcommand{\paragraph}[1]{\oldparagraph{#1}\mbox{}}
\fi
\ifx\subparagraph\undefined\else
  \let\oldsubparagraph\subparagraph
  \renewcommand{\subparagraph}[1]{\oldsubparagraph{#1}\mbox{}}
\fi

% Lower the position of page numbers
\usepackage{fancyhdr}
\fancyhf{}
\fancyfoot[C]{\thepage}
\setlength{\footskip}{100pt} % Increase this value to lower the page numbers


\usepackage{color}
\usepackage{fancyvrb}
\newcommand{\VerbBar}{|}
\newcommand{\VERB}{\Verb[commandchars=\\\{\}]}
\DefineVerbatimEnvironment{Highlighting}{Verbatim}{commandchars=\\\{\}}
% Add ',fontsize=\small' for more characters per line
\usepackage{framed}
\definecolor{shadecolor}{RGB}{241,243,245}
\newenvironment{Shaded}{\begin{snugshade}}{\end{snugshade}}
\newcommand{\AlertTok}[1]{\textcolor[rgb]{0.68,0.00,0.00}{#1}}
\newcommand{\AnnotationTok}[1]{\textcolor[rgb]{0.37,0.37,0.37}{#1}}
\newcommand{\AttributeTok}[1]{\textcolor[rgb]{0.40,0.45,0.13}{#1}}
\newcommand{\BaseNTok}[1]{\textcolor[rgb]{0.68,0.00,0.00}{#1}}
\newcommand{\BuiltInTok}[1]{\textcolor[rgb]{0.00,0.23,0.31}{#1}}
\newcommand{\CharTok}[1]{\textcolor[rgb]{0.13,0.47,0.30}{#1}}
\newcommand{\CommentTok}[1]{\textcolor[rgb]{0.37,0.37,0.37}{#1}}
\newcommand{\CommentVarTok}[1]{\textcolor[rgb]{0.37,0.37,0.37}{\textit{#1}}}
\newcommand{\ConstantTok}[1]{\textcolor[rgb]{0.56,0.35,0.01}{#1}}
\newcommand{\ControlFlowTok}[1]{\textcolor[rgb]{0.00,0.23,0.31}{\textbf{#1}}}
\newcommand{\DataTypeTok}[1]{\textcolor[rgb]{0.68,0.00,0.00}{#1}}
\newcommand{\DecValTok}[1]{\textcolor[rgb]{0.68,0.00,0.00}{#1}}
\newcommand{\DocumentationTok}[1]{\textcolor[rgb]{0.37,0.37,0.37}{\textit{#1}}}
\newcommand{\ErrorTok}[1]{\textcolor[rgb]{0.68,0.00,0.00}{#1}}
\newcommand{\ExtensionTok}[1]{\textcolor[rgb]{0.00,0.23,0.31}{#1}}
\newcommand{\FloatTok}[1]{\textcolor[rgb]{0.68,0.00,0.00}{#1}}
\newcommand{\FunctionTok}[1]{\textcolor[rgb]{0.28,0.35,0.67}{#1}}
\newcommand{\ImportTok}[1]{\textcolor[rgb]{0.00,0.46,0.62}{#1}}
\newcommand{\InformationTok}[1]{\textcolor[rgb]{0.37,0.37,0.37}{#1}}
\newcommand{\KeywordTok}[1]{\textcolor[rgb]{0.00,0.23,0.31}{\textbf{#1}}}
\newcommand{\NormalTok}[1]{\textcolor[rgb]{0.00,0.23,0.31}{#1}}
\newcommand{\OperatorTok}[1]{\textcolor[rgb]{0.37,0.37,0.37}{#1}}
\newcommand{\OtherTok}[1]{\textcolor[rgb]{0.00,0.23,0.31}{#1}}
\newcommand{\PreprocessorTok}[1]{\textcolor[rgb]{0.68,0.00,0.00}{#1}}
\newcommand{\RegionMarkerTok}[1]{\textcolor[rgb]{0.00,0.23,0.31}{#1}}
\newcommand{\SpecialCharTok}[1]{\textcolor[rgb]{0.37,0.37,0.37}{#1}}
\newcommand{\SpecialStringTok}[1]{\textcolor[rgb]{0.13,0.47,0.30}{#1}}
\newcommand{\StringTok}[1]{\textcolor[rgb]{0.13,0.47,0.30}{#1}}
\newcommand{\VariableTok}[1]{\textcolor[rgb]{0.07,0.07,0.07}{#1}}
\newcommand{\VerbatimStringTok}[1]{\textcolor[rgb]{0.13,0.47,0.30}{#1}}
\newcommand{\WarningTok}[1]{\textcolor[rgb]{0.37,0.37,0.37}{\textit{#1}}}

\providecommand{\tightlist}{%
  \setlength{\itemsep}{0pt}\setlength{\parskip}{0pt}}\usepackage{longtable,booktabs,array}
\usepackage{calc} % for calculating minipage widths
% Correct order of tables after \paragraph or \subparagraph
\usepackage{etoolbox}
\makeatletter
\patchcmd\longtable{\par}{\if@noskipsec\mbox{}\fi\par}{}{}
\makeatother
% Allow footnotes in longtable head/foot
\IfFileExists{footnotehyper.sty}{\usepackage{footnotehyper}}{\usepackage{footnote}}
\makesavenoteenv{longtable}
\usepackage{graphicx}
\makeatletter
\newsavebox\pandoc@box
\newcommand*\pandocbounded[1]{% scales image to fit in text height/width
  \sbox\pandoc@box{#1}%
  \Gscale@div\@tempa{\textheight}{\dimexpr\ht\pandoc@box+\dp\pandoc@box\relax}%
  \Gscale@div\@tempb{\linewidth}{\wd\pandoc@box}%
  \ifdim\@tempb\p@<\@tempa\p@\let\@tempa\@tempb\fi% select the smaller of both
  \ifdim\@tempa\p@<\p@\scalebox{\@tempa}{\usebox\pandoc@box}%
  \else\usebox{\pandoc@box}%
  \fi%
}
% Set default figure placement to htbp
\def\fps@figure{htbp}
\makeatother

% Change font size
\usepackage[font={footnotesize, sf}, labelfont=bf, margin=0.05\linewidth]{caption}
\renewcommand{\familydefault}{\sfdefault}
% Bold the 'Figure #' in the caption and separate it from the title/caption
% with a period
% Captions will be left justified
\usepackage[
	aboveskip=1pt,
	labelfont=bf,
	labelsep=period,
	justification=raggedright,
	singlelinecheck=off
]{caption}
\renewcommand{\figurename}{Fig}

% Add numbered lines
\usepackage{lineno}
% \linenumbers

% Package to include multiple title pages
% This will allow me to add a tile to the main text and to the SI
\usepackage{titling}

% Define command to begin the supplementary section
\newcommand{\beginsupplement}{
\setcounter{page}{1} % Restart page counter
\setcounter{section}{0} % Restart section counter
% \renewcommand{\thesection}{\Alph{section}}%
\renewcommand{\thesection}{}
\renewcommand{\thesubsection}{\Alph{subsection}}
\setcounter{table}{0} % Restart table counter
\renewcommand{\thetable}{\Alph{table}}
\setcounter{figure}{0} % Restart figure counter
\renewcommand{\thefigure}{\Alph{figure}}%
\setcounter{equation}{0} % Restart equation counter
\renewcommand{\theequation}{S\arabic{equation}}%
}

% Add author affiliations
\usepackage{authblk}

% Allow to define personalized colors
\usepackage{xcolor}
\makeatletter
\@ifpackageloaded{caption}{}{\usepackage{caption}}
\AtBeginDocument{%
\ifdefined\contentsname
  \renewcommand*\contentsname{Table of contents}
\else
  \newcommand\contentsname{Table of contents}
\fi
\ifdefined\listfigurename
  \renewcommand*\listfigurename{List of Figures}
\else
  \newcommand\listfigurename{List of Figures}
\fi
\ifdefined\listtablename
  \renewcommand*\listtablename{List of Tables}
\else
  \newcommand\listtablename{List of Tables}
\fi
\ifdefined\figurename
  \renewcommand*\figurename{Figure}
\else
  \newcommand\figurename{Figure}
\fi
\ifdefined\tablename
  \renewcommand*\tablename{Table}
\else
  \newcommand\tablename{Table}
\fi
}
\@ifpackageloaded{float}{}{\usepackage{float}}
\floatstyle{ruled}
\@ifundefined{c@chapter}{\newfloat{codelisting}{h}{lop}}{\newfloat{codelisting}{h}{lop}[chapter]}
\floatname{codelisting}{Listing}
\newcommand*\listoflistings{\listof{codelisting}{List of Listings}}
\makeatother
\makeatletter
\makeatother
\makeatletter
\@ifpackageloaded{caption}{}{\usepackage{caption}}
\@ifpackageloaded{subcaption}{}{\usepackage{subcaption}}
\makeatother
\ifLuaTeX
  \usepackage{selnolig}  % disable illegal ligatures
\fi

%%%%%%%%%%%%%%%%%%%%%%%%%%%%%%%%%%%%%%%%%%%%%%%%%%%%%%%%%%%%%%%%%%%%%%%%%%%%%%%%
\usepackage[
    style=vancouver,
    sorting=none,				% Do not sort bibliography
	url=false, 					% Do not show url in reference
	doi=false, 					% Do not show doi in reference
	isbn=false, 				% Do not show isbn link in reference
	eprint=false, 			    % Do not show eprint link in reference
	maxbibnames=9, 			    % Include up to 9 names in citation
	firstinits=true,
]{biblatex}
%%%%%%%%%%%%%%%%%%%%%%%%%%%%%%%%%%%%%%%%%%%%%%%%%%%%%%%%%%%%%%%%%%%%%%%%%%%%%%%%
\addbibresource{../references.bib}
\IfFileExists{bookmark.sty}{\usepackage{bookmark}}{\usepackage{hyperref}}
\IfFileExists{xurl.sty}{\usepackage{xurl}}{} % add URL line breaks if available
\urlstyle{same} % disable monospaced font for URLs
\hypersetup{
  pdftitle={Bayesian Inference of IC\_\{50\} Values},
  pdfauthor={Manuel Razo-Mejia},
  pdfkeywords={Bayesian inference, Antibiotic resistance, \(IC_{50}\)},
  colorlinks=true,
  linkcolor={blue},
  filecolor={Maroon},
  citecolor={Blue},
  urlcolor={Blue},
  pdfcreator={LaTeX via pandoc}}

%%%%%%%%%%%%%%%%%%%%%%%%%%%%%%%%%%%%%%%%%%%%%%%%%%%%%%%%%%%%%%%%%%%%%%%%%%%%%%%%
% Add title
\title{Bayesian Inference of \(IC_{50}\) Values}
%%%%%%%%%%%%%%%%%%%%%%%%%%%%%%%%%%%%%%%%%%%%%%%%%%%%%%%%%%%%%%%%%%%%%%%%%%%%%%%%

%%%%%%%%%%%%%%%%%%%%%%%%%%%%%%%%%%%%%%%%%%%%%%%%%%%%%%%%%%%%%%%%%%%%%%%%%%%%%%%%
% Add list of authors

% Loop through authors
\author[1]{Manuel Razo-Mejia}

% Loop through affiliations
\affil[1]{Department of Biology, Stanford University}

% Add corresponding author
%%%%%%%%%%%%%%%%%%%%%%%%%%%%%%%%%%%%%%%%%%%%%%%%%%%%%%%%%%%%%%%%%%%%%%%%%%%%%%%%

%%%%%%%%%%%%%%%%%%%%%%%%%%%%%%%%%%%%%%%%%%%%%%%%%%%%%%%%%%%%%%%%%%%%%%%%%%%%%%%%
% Set affiliations in small font
\renewcommand\Affilfont{\itshape\small}
%%%%%%%%%%%%%%%%%%%%%%%%%%%%%%%%%%%%%%%%%%%%%%%%%%%%%%%%%%%%%%%%%%%%%%%%%%%%%%%%

%%%%%%%%%%%%%%%%%%%%%%%%%%%%%%%%%%%%%%%%%%%%%%%%%%%%%%%%%%%%%%%%%%%%%%%%%%%%%%%%
% Remove date
\date{}
%%%%%%%%%%%%%%%%%%%%%%%%%%%%%%%%%%%%%%%%%%%%%%%%%%%%%%%%%%%%%%%%%%%%%%%%%%%%%%%%

% Add package for line numbering
\usepackage{lineno}

\begin{document}
\linenumbers
\maketitle
(c) This work is licensed under a
\href{https://creativecommons.org/licenses/by/4.0/}{Creative Commons
Attribution License CC-BY 4.0}. All code contained herein is licensed
under an \href{https://opensource.org/licenses/MIT}{MIT license}.

\begin{Shaded}
\begin{Highlighting}[]
\CommentTok{\# Import project package}
\ImportTok{import} \BuiltInTok{Antibiotic}

\CommentTok{\# Import package to handle DataFrames}
\ImportTok{import} \BuiltInTok{DataFrames}\NormalTok{ as DF}
\ImportTok{import} \BuiltInTok{CSV}

\CommentTok{\# Import library for Bayesian inference}
\ImportTok{import} \BuiltInTok{Turing}

\CommentTok{\# Import library to list files}
\ImportTok{import} \BuiltInTok{Glob}

\CommentTok{\# Import packages to work with data}
\ImportTok{import} \BuiltInTok{DataFrames}\NormalTok{ as DF}

\CommentTok{\# Load CairoMakie for plotting}
\ImportTok{using} \BuiltInTok{CairoMakie}
\ImportTok{import} \BuiltInTok{ColorSchemes}

\CommentTok{\# Import packages for posterior sampling}
\ImportTok{import} \BuiltInTok{PairPlots}

\CommentTok{\# Import basic math libraries}
\ImportTok{import} \BuiltInTok{LsqFit}
\ImportTok{import} \BuiltInTok{StatsBase}
\ImportTok{import} \BuiltInTok{LinearAlgebra}
\ImportTok{import} \BuiltInTok{Random}

\CommentTok{\# Activate backend}
\NormalTok{CairoMakie.}\FunctionTok{activate!}\NormalTok{()}

\CommentTok{\# Set custom plotting style}
\NormalTok{Antibiotic.viz.}\FunctionTok{theme\_makie!}\NormalTok{()}
\end{Highlighting}
\end{Shaded}

(c) This work is licensed under a
\href{https://creativecommons.org/licenses/by/4.0/}{Creative Commons
Attribution License CC-BY 4.0}. All code contained herein is licensed
under an \href{https://opensource.org/licenses/MIT}{MIT license}.

\section{\texorpdfstring{Bayesian Inference of \(IC_{50}\)
values}{Bayesian Inference of IC\_\{50\} values}}\label{bayesian-inference-of-ic_50-values}

In this notebook, we will perform Bayesian inference on the \(IC_{50}\)
values of the antibiotic resistance landscape. For this, we will use the
raw \texttt{OD620} measurements provided by \textcite{iwasawa2022}.

Let's begin by loading the data into a DataFrame.

\begin{Shaded}
\begin{Highlighting}[]
\CommentTok{\# Load data into a DataFrame}
\NormalTok{df }\OperatorTok{=}\NormalTok{ CSV.}\FunctionTok{read}\NormalTok{(}\StringTok{"./iwasawa\_data/iwasawa\_tidy.csv"}\NormalTok{, DF.DataFrame)}

\FunctionTok{first}\NormalTok{(df, }\FloatTok{5}\NormalTok{)}
\end{Highlighting}
\end{Shaded}

\begin{tabular}{r|ccccccccc}
    & antibiotic & col & OD & design & concentration\_ugmL & day & strain\_num & env & \\
    \hline
    & String7 & String7 & Float64 & Int64 & Float64 & Int64 & Int64 & String15 & \\
    \hline
    1 & TET & col1 & 0.0464 & 3 & 0.0 & 1 & 19 & TETE4\_in\_KM & $\dots$ \\
    2 & KM & col1 & 0.0432 & 3 & 0.0 & 1 & 19 & TETE4\_in\_KM & $\dots$ \\
    3 & NFLX & col1 & 0.0414 & 3 & 0.0 & 1 & 19 & TETE4\_in\_KM & $\dots$ \\
    4 & SS & col1 & 0.0415 & 3 & 0.0 & 1 & 19 & TETE4\_in\_KM & $\dots$ \\
    5 & PLM & col1 & 0.0427 & 3 & 0.0 & 1 & 19 & TETE4\_in\_KM & $\dots$ \\
\end{tabular}

To double-check that the structure of the table makes sense, let's plot
the time series for one example to see if the sequence agrees with the
expectation.

\begin{Shaded}
\begin{Highlighting}[]
\CommentTok{\# Define data to use}
\NormalTok{data }\OperatorTok{=}\NormalTok{ df[}
\NormalTok{    (df.antibiotic}\OperatorTok{.==}\StringTok{"KM"}\NormalTok{)}\OperatorTok{.\&}\NormalTok{(df.env}\OperatorTok{.==}\StringTok{"Parent\_in\_KM"}\NormalTok{)}\OperatorTok{.\&}\NormalTok{(df.strain\_num}\OperatorTok{.==}\FloatTok{13}\NormalTok{)}\OperatorTok{.\&}\NormalTok{.!(df.blank)}\OperatorTok{.\&}\NormalTok{(df.concentration\_ugmL}\OperatorTok{.\textgreater{}}\FloatTok{0}\NormalTok{),}
    \OperatorTok{:}\NormalTok{]}
\CommentTok{\# Remove blank measurement}
\CommentTok{\# Group data by day}
\NormalTok{df\_group }\OperatorTok{=}\NormalTok{ DF.}\FunctionTok{groupby}\NormalTok{(data, }\OperatorTok{:}\NormalTok{day)}

\CommentTok{\# Initialize figure}
\NormalTok{fig }\OperatorTok{=} \FunctionTok{Figure}\NormalTok{(size}\OperatorTok{=}\NormalTok{(}\FloatTok{500}\NormalTok{, }\FloatTok{300}\NormalTok{))}

\CommentTok{\# Add axis}
\NormalTok{ax }\OperatorTok{=} \FunctionTok{Axis}\NormalTok{(}
\NormalTok{    fig[}\FloatTok{1}\NormalTok{, }\FloatTok{1}\NormalTok{],}
\NormalTok{    xlabel}\OperatorTok{=}\StringTok{"antibiotic concentration"}\NormalTok{,}
\NormalTok{    ylabel}\OperatorTok{=}\StringTok{"OD₆₂₀"}\NormalTok{,}
\NormalTok{    xscale}\OperatorTok{=}\NormalTok{log2}
\NormalTok{)}

\CommentTok{\# Define colors for plot}
\NormalTok{colors }\OperatorTok{=} \FunctionTok{get}\NormalTok{(ColorSchemes.Blues\_9, }\FunctionTok{LinRange}\NormalTok{(}\FloatTok{0.25}\NormalTok{, }\FloatTok{1}\NormalTok{, }\FunctionTok{length}\NormalTok{(df\_group)))}

\CommentTok{\# Loop through days}
\ControlFlowTok{for}\NormalTok{ (i, d) }\KeywordTok{in} \FunctionTok{enumerate}\NormalTok{(df\_group)}
    \CommentTok{\# Sort data by concentration}
\NormalTok{    DF.}\FunctionTok{sort!}\NormalTok{(d, }\OperatorTok{:}\NormalTok{concentration\_ugmL)}
    \CommentTok{\# Plot scatter line}
    \FunctionTok{scatterlines!}\NormalTok{(}
\NormalTok{        ax, d.concentration\_ugmL, d.OD, color}\OperatorTok{=}\NormalTok{colors[i], label}\OperatorTok{=}\StringTok{"}\SpecialCharTok{$}\NormalTok{(}\FunctionTok{first}\NormalTok{(d.day))}\StringTok{"}
\NormalTok{    )}
\ControlFlowTok{end} \CommentTok{\# for}

\CommentTok{\# Add legend to plot}
\NormalTok{fig[}\FloatTok{1}\NormalTok{, }\FloatTok{2}\NormalTok{] }\OperatorTok{=} \FunctionTok{Legend}\NormalTok{(}
\NormalTok{    fig, ax, }\StringTok{"day"}\NormalTok{, framevisible}\OperatorTok{=}\ConstantTok{false}\NormalTok{, nbanks}\OperatorTok{=}\FloatTok{3}\NormalTok{, labelsize}\OperatorTok{=}\FloatTok{10}
\NormalTok{)}

\NormalTok{fig}
\end{Highlighting}
\end{Shaded}

\includegraphics[width=10.41667in,height=6.25in]{iwasawa_bayesian_inference_files/figure-pdf/cell-4-output-1.png}

The functional form used by the authors to fit the data is
\begin{equation}\phantomsection\label{eq-iwasawa-model}{
f(x) = \frac{a}
{1+\exp \left[b\left(\log _2 x-\log _2 IC_{50}\right)\right]} + c
}\end{equation}

where \(a\), \(b\), and \(c\) are nuisance parameters of the model,
\(IC_{50}\) is the parameter of interest, and \(x\) is the antibiotic
concentration. We can define a function to compute this model.

\begin{Shaded}
\begin{Highlighting}[]
\PreprocessorTok{@doc} \StringTok{raw"""}
\StringTok{    logistic(x, a, b, c, ic50)}

\StringTok{Compute the logistic function used to model the relationship between antibiotic}
\StringTok{concentration and bacterial growth.}

\StringTok{This function implements the following equation:}

\StringTok{f(x) = a / (1 + exp(b * (log₂(x) {-} log₂(IC₅₀)))) + c}

\StringTok{\# Arguments}
\StringTok{{-} \textasciigrave{}x\textasciigrave{}: Antibiotic concentration (input variable)}
\StringTok{{-} \textasciigrave{}a\textasciigrave{}: Maximum effect parameter (difference between upper and lower asymptotes)}
\StringTok{{-} \textasciigrave{}b\textasciigrave{}: Slope parameter (steepness of the curve)}
\StringTok{{-} \textasciigrave{}c\textasciigrave{}: Minimum effect parameter (lower asymptote)}
\StringTok{{-} \textasciigrave{}ic50\textasciigrave{}: IC₅₀ parameter (concentration at which the effect is halfway between}
\StringTok{  the minimum and maximum)}

\StringTok{\# Returns}
\StringTok{The computed effect (e.g., optical density) for the given antibiotic}
\StringTok{concentration and parameters.}

\StringTok{Note: This function is vectorized and can handle array inputs for \textasciigrave{}x\textasciigrave{}.}
\StringTok{"""}
\KeywordTok{function} \FunctionTok{logistic}\NormalTok{(x, a, b, c, ic50)}
    \ControlFlowTok{return}\NormalTok{ @. a }\OperatorTok{/}\NormalTok{ (}\FloatTok{1.0} \OperatorTok{+} \FunctionTok{exp}\NormalTok{(b }\OperatorTok{*}\NormalTok{ (}\FunctionTok{log2}\NormalTok{(x) }\OperatorTok{{-}} \FunctionTok{log2}\NormalTok{(ic50)))) }\OperatorTok{+}\NormalTok{ c}
\KeywordTok{end}
\end{Highlighting}
\end{Shaded}

To test the function, let's plot the model for a set of parameters.

\begin{Shaded}
\begin{Highlighting}[]
\CommentTok{\# Define parameters}
\NormalTok{a }\OperatorTok{=} \FloatTok{1.0}
\NormalTok{b }\OperatorTok{=} \FloatTok{1.0}
\NormalTok{c }\OperatorTok{=} \FloatTok{0.0}
\NormalTok{ic50 }\OperatorTok{=} \FloatTok{1.0}

\CommentTok{\# Define concentration range}
\NormalTok{x }\OperatorTok{=} \FloatTok{10} \OperatorTok{.\^{}} \FunctionTok{LinRange}\NormalTok{(}\OperatorTok{{-}}\FloatTok{2.5}\NormalTok{, }\FloatTok{2.5}\NormalTok{, }\FloatTok{50}\NormalTok{)}

\CommentTok{\# Compute model}
\NormalTok{y }\OperatorTok{=} \FunctionTok{logistic}\NormalTok{(x, a, b, c, ic50)}

\CommentTok{\# Initialize figure}
\NormalTok{fig }\OperatorTok{=} \FunctionTok{Figure}\NormalTok{(size}\OperatorTok{=}\NormalTok{(}\FloatTok{350}\NormalTok{, }\FloatTok{300}\NormalTok{))}
\CommentTok{\# Add axis}
\NormalTok{ax }\OperatorTok{=} \FunctionTok{Axis}\NormalTok{(}
\NormalTok{    fig[}\FloatTok{1}\NormalTok{, }\FloatTok{1}\NormalTok{],}
\NormalTok{    xlabel}\OperatorTok{=}\StringTok{"antibiotic concentration (a.u.)"}\NormalTok{,}
\NormalTok{    ylabel}\OperatorTok{=}\StringTok{"optical density"}\NormalTok{,}
\NormalTok{    xscale}\OperatorTok{=}\NormalTok{log10}
\NormalTok{)}
\CommentTok{\# Plot model}
\FunctionTok{lines!}\NormalTok{(ax, x, y)}

\NormalTok{fig}
\end{Highlighting}
\end{Shaded}

\includegraphics[width=7.29167in,height=6.25in]{iwasawa_bayesian_inference_files/figure-pdf/cell-6-output-1.png}

The function seems to work as expected. However, notice that in Eq. (1),
the \(\mathrm{IC}_{50}\) goes into the logarithm. This is the parameter
we will fit for. Thus, let's define a function that takes
\(\log_2(\mathrm{IC}_{50})\) and \(\log_2(x)\) as input instead.

\begin{Shaded}
\begin{Highlighting}[]
\PreprocessorTok{@doc} \StringTok{raw"""}
\StringTok{    logistic\_log2(log2x, a, b, c, log2ic50)}

\StringTok{Compute the logistic function used to model the relationship between antibiotic}
\StringTok{concentration and bacterial growth, using log2 inputs for concentration and}
\StringTok{IC₅₀.}

\StringTok{This function implements the following equation:}

\StringTok{f(x) = a / (1 + exp(b * (log₂(x) {-} log₂(IC₅₀)))) + c}

\StringTok{\# Arguments}
\StringTok{{-} \textasciigrave{}log2x\textasciigrave{}: log₂ of the antibiotic concentration (input variable)}
\StringTok{{-} \textasciigrave{}a\textasciigrave{}: Maximum effect parameter (difference between upper and lower asymptotes)}
\StringTok{{-} \textasciigrave{}b\textasciigrave{}: Slope parameter (steepness of the curve)}
\StringTok{{-} \textasciigrave{}c\textasciigrave{}: Minimum effect parameter (lower asymptote)}
\StringTok{{-} \textasciigrave{}log2ic50\textasciigrave{}: log₂ of the IC₅₀ parameter}

\StringTok{\# Returns}
\StringTok{The computed effect (e.g., optical density) for the given log₂ antibiotic}
\StringTok{concentration and parameters.}

\StringTok{Note: This function is vectorized and can handle array inputs for \textasciigrave{}log2x\textasciigrave{}.}
\StringTok{"""}
\KeywordTok{function} \FunctionTok{logistic\_log2}\NormalTok{(log2x, a, b, c, log2ic50)}
    \ControlFlowTok{return}\NormalTok{ @. a }\OperatorTok{/}\NormalTok{ (}\FloatTok{1.0} \OperatorTok{+} \FunctionTok{exp}\NormalTok{(b }\OperatorTok{*}\NormalTok{ (log2x }\OperatorTok{{-}}\NormalTok{ log2ic50))) }\OperatorTok{+}\NormalTok{ c}
\KeywordTok{end}
\end{Highlighting}
\end{Shaded}

\subsection{Bayesian model}\label{bayesian-model}

Given the model presented in Eq. (1), and the data, our objective is to
infer the value of all parameters. By Bayes theorem, we write

\begin{equation}\phantomsection\label{eq-iwasawa-bayesian-model}{
\pi(IC_{50}, a, b, c \mid \text{data}) = 
\frac{\pi(\text{data} \mid IC_{50}, a, b, c) 
\pi(IC_{50}, a, b, c)}
{\pi(\text{data})},
}\end{equation}

where \(\text{data}\) consists of the pairs of antibiotic concentration
and optical density.

\subsubsection{\texorpdfstring{Likelihood
\(\pi(\text{data} \mid IC_{50}, a, b, c)\)}{Likelihood \textbackslash pi(\textbackslash text\{data\} \textbackslash mid IC\_\{50\}, a, b, c)}}\label{likelihood-pitextdata-mid-ic_50-a-b-c}

Let's begin by defining the likelihood function. For simplicity, we
assume each datum is independent and identically distributed (i.i.d.)
and write

\begin{equation}\phantomsection\label{eq-iwasawa-likelihood}{
\pi(\text{data} \mid IC_{50}, a, b, c) = 
\prod_{i=1}^n \pi(d_i \mid IC_{50}, a, b, c),
}\end{equation}

where \(d_i = (x_i, y_i)\) is the \(i\)-th pair of antibiotic
concentration and optical density, respectively, and \(n\) is the total
number of data points. As a first pass, we assume that our experimental
measurements can be expressed as

\begin{equation}\phantomsection\label{eq-iwasawa-likelihood-model}{
y_i = f(x_i, IC_{50}, a, b, c) + \epsilon_i,
}\end{equation}

where \(\epsilon_i\) is the experimental error. Furthermore, we assume
that the experimental error is normally distributed, i.e.,

\begin{equation}\phantomsection\label{eq-iwasawa-likelihood-error}{
\epsilon_i \sim \mathcal{N}(0, \sigma^2),
}\end{equation}

where \(\sigma^2\) is an unknown variance parameter that must be
included in our inference. Notice that we assume the same variance
parameter for all data points since \(\sigma^2\) is not indexed by
\(i\).

Given this likelihood function, we must update our inference on the
parameters as

\begin{equation}\phantomsection\label{eq-iwasawa-likelihood-error-model}{
\pi(IC_{50}, a, b, c, \sigma^2 \mid \text{data}) = 
\frac{\pi(\text{data} \mid IC_{50}, a, b, c, \sigma^2) 
\pi(IC_{50}, a, b, c, \sigma^2)}
{\pi(\text{data})},
}\end{equation}

to include the new parameter \(\sigma^2\).

Our likelihood function is then of the form
\begin{equation}\phantomsection\label{eq-iwasawa-likelihood-error-model}{
y_i \mid IC_{50}, a, b, c, \sigma^2 \sim 
\mathcal{N}(f(x_i, IC_{50}, a, b, c), \sigma^2).
}\end{equation}

\subsubsection{\texorpdfstring{Prior
\(\pi(IC_{50}, a, b, c, \sigma^2)\)}{Prior \textbackslash pi(IC\_\{50\}, a, b, c, \textbackslash sigma\^{}2)}}\label{prior-piic_50-a-b-c-sigma2}

For the prior, we assume that all parameters are independent and write
\begin{equation}\phantomsection\label{eq-iwasawa-prior}{
\pi(IC_{50}, a, b, c, \sigma^2) = 
\pi(IC_{50}) \pi(a) \pi(b) \pi(c) \pi(\sigma^2).
}\end{equation}

Let's detail each prior.

\begin{enumerate}
\def\labelenumi{\arabic{enumi}.}
\tightlist
\item
  \(IC_{50}\): The \(IC_{50}\) is a strictly positive parameter.
  However, we will fit for \(\log_2(IC_{50})\). Thus, we will use a
  normal prior for \(\log_2(IC_{50})\). This means we have
\end{enumerate}

\begin{equation}\phantomsection\label{eq-iwasawa-prior-ic50}{
\log_2(IC_{50}) \sim 
\mathcal{N}(
    \mu_{\log_2(IC_{50})}, \sigma_{\log_2(IC_{50})}^2
).
}\end{equation}

\begin{enumerate}
\def\labelenumi{\arabic{enumi}.}
\setcounter{enumi}{1}
\tightlist
\item
  \(a\): This nuisance parameter scales the logistic function. Again,
  the natural scale for this parameter is a strictly positive real
  number. Thus, we will use a lognormal prior for \(a\). This means we
  have
\end{enumerate}

\begin{equation}\phantomsection\label{eq-iwasawa-prior-a}{
a \sim \text{LogNormal}(\mu_a, \sigma_a^2).
}\end{equation}

\begin{enumerate}
\def\labelenumi{\arabic{enumi}.}
\setcounter{enumi}{2}
\tightlist
\item
  \(b\): This parameter controls the steepness of the logistic function.
  Again, the natural scale for this parameter is a strictly positive
  real number. Thus, we will use a lognormal prior for \(b\). This means
  we have
\end{enumerate}

\begin{equation}\phantomsection\label{eq-iwasawa-prior-b}{
b \sim \text{LogNormal}(\mu_b, \sigma_b^2).
}\end{equation}

\begin{enumerate}
\def\labelenumi{\arabic{enumi}.}
\setcounter{enumi}{3}
\tightlist
\item
  \(c\): This parameter controls the minimum value of the logistic
  function. Since this is a strictly positive real number that does not
  necessarily scale with the data, we will use a half-normal prior for
  \(c\). This means we have
\end{enumerate}

\begin{equation}\phantomsection\label{eq-iwasawa-prior-c}{
c \sim \text{Half-}\mathcal{N}(0, \sigma_c^2).
}\end{equation}

\begin{enumerate}
\def\labelenumi{\arabic{enumi}.}
\setcounter{enumi}{4}
\tightlist
\item
  \(\sigma^2\): This parameter controls the variance of the experimental
  error. Since this is a strictly positive real number that does not
  necessarily scale with the data, we will use a half-normal prior for
  \(\sigma^2\). This means we have
\end{enumerate}

\begin{equation}\phantomsection\label{eq-iwasawa-prior-sigma2}{
\sigma^2 \sim \text{Half-}\mathcal{N}(0, \sigma_{\sigma^2}^2).
}\end{equation}

With all of this in place, we are ready to define a \texttt{Turing}
model to perform Bayesian inference on the parameters of the model.

\begin{Shaded}
\begin{Highlighting}[]
\NormalTok{Turing.}\PreprocessorTok{@model} \KeywordTok{function} \FunctionTok{logistic\_model}\NormalTok{(}
\NormalTok{    log2x, y, prior\_params}\OperatorTok{::}\DataTypeTok{NamedTuple}\NormalTok{=}\FunctionTok{NamedTuple}\NormalTok{()}
\NormalTok{)}
    \CommentTok{\# Define default prior parameters}
\NormalTok{    default\_params }\OperatorTok{=}\NormalTok{ (}
\NormalTok{        log2ic50}\OperatorTok{=}\NormalTok{(}\FloatTok{0}\NormalTok{, }\FloatTok{1}\NormalTok{),}
\NormalTok{        a}\OperatorTok{=}\NormalTok{(}\FloatTok{0}\NormalTok{, }\FloatTok{1}\NormalTok{),}
\NormalTok{        b}\OperatorTok{=}\NormalTok{(}\FloatTok{0}\NormalTok{, }\FloatTok{1}\NormalTok{),}
\NormalTok{        c}\OperatorTok{=}\NormalTok{(}\FloatTok{0}\NormalTok{, }\FloatTok{1}\NormalTok{),}
\NormalTok{        σ²}\OperatorTok{=}\NormalTok{(}\FloatTok{0}\NormalTok{, }\FloatTok{1}\NormalTok{)}
\NormalTok{    )}

    \CommentTok{\# Merge default parameters with provided parameters}
\NormalTok{    params }\OperatorTok{=} \FunctionTok{merge}\NormalTok{(default\_params, prior\_params)}

    \CommentTok{\# Define priors}
\NormalTok{    log2ic50 }\OperatorTok{\textasciitilde{}}\NormalTok{ Turing.}\FunctionTok{Normal}\NormalTok{(params.log2ic50}\OperatorTok{...}\NormalTok{)}
\NormalTok{    a }\OperatorTok{\textasciitilde{}}\NormalTok{ Turing.}\FunctionTok{LogNormal}\NormalTok{(params.a}\OperatorTok{...}\NormalTok{)}
\NormalTok{    b }\OperatorTok{\textasciitilde{}}\NormalTok{ Turing.}\FunctionTok{LogNormal}\NormalTok{(params.b}\OperatorTok{...}\NormalTok{)}
\NormalTok{    c }\OperatorTok{\textasciitilde{}}\NormalTok{ Turing.}\FunctionTok{truncated}\NormalTok{(Turing.}\FunctionTok{Normal}\NormalTok{(params.c}\OperatorTok{...}\NormalTok{), }\FloatTok{0}\NormalTok{, }\ConstantTok{Inf}\NormalTok{)}
\NormalTok{    σ² }\OperatorTok{\textasciitilde{}}\NormalTok{ Turing.}\FunctionTok{truncated}\NormalTok{(Turing.}\FunctionTok{Normal}\NormalTok{(params.σ²}\OperatorTok{...}\NormalTok{), }\FloatTok{0}\NormalTok{, }\ConstantTok{Inf}\NormalTok{)}

    \CommentTok{\# Define likelihood}
\NormalTok{    y }\OperatorTok{\textasciitilde{}}\NormalTok{ Turing.}\FunctionTok{MvNormal}\NormalTok{(}
        \FunctionTok{logistic\_log2}\NormalTok{(log2x, a, b, c, log2ic50),}
        \BuiltInTok{LinearAlgebra}\NormalTok{.}\FunctionTok{Diagonal}\NormalTok{(}\FunctionTok{fill}\NormalTok{(σ², }\FunctionTok{length}\NormalTok{(y)))}
\NormalTok{    )}
\KeywordTok{end}
\end{Highlighting}
\end{Shaded}

Having defined the model, we can now perform inference. First, let's
perform inference on simulated data. For this, we first simulate a
single titration curve with known parameters.

\begin{Shaded}
\begin{Highlighting}[]
\BuiltInTok{Random}\NormalTok{.}\FunctionTok{seed!}\NormalTok{(}\FloatTok{42}\NormalTok{)}
\CommentTok{\# Define ground truth parameters}
\NormalTok{log2ic50\_true }\OperatorTok{=} \FloatTok{0.5}
\NormalTok{a\_true }\OperatorTok{=} \FloatTok{1.0}
\NormalTok{b\_true }\OperatorTok{=} \FloatTok{10.0}
\NormalTok{σ²\_true }\OperatorTok{=} \FloatTok{0.01}
\NormalTok{c\_true }\OperatorTok{=} \FloatTok{0.1}

\CommentTok{\# Simulate data}
\NormalTok{log2x }\OperatorTok{=} \FunctionTok{LinRange}\NormalTok{(}\OperatorTok{{-}}\FloatTok{2.5}\NormalTok{, }\FloatTok{2.5}\NormalTok{, }\FloatTok{15}\NormalTok{)}
\CommentTok{\# Define mean of data}
\NormalTok{ŷ }\OperatorTok{=} \FunctionTok{logistic\_log2}\NormalTok{(log2x, a\_true, b\_true, c\_true, log2ic50\_true)}
\CommentTok{\# Add noise}
\NormalTok{y }\OperatorTok{=}\NormalTok{ ŷ }\OperatorTok{.+} \FunctionTok{randn}\NormalTok{(}\FunctionTok{length}\NormalTok{(log2x)) }\OperatorTok{*} \FunctionTok{√}\NormalTok{(σ²\_true)}

\CommentTok{\# Initialize figure}
\NormalTok{fig }\OperatorTok{=} \FunctionTok{Figure}\NormalTok{(size}\OperatorTok{=}\NormalTok{(}\FloatTok{350}\NormalTok{, }\FloatTok{300}\NormalTok{))}
\CommentTok{\# Add axis}
\NormalTok{ax }\OperatorTok{=} \FunctionTok{Axis}\NormalTok{(}
\NormalTok{    fig[}\FloatTok{1}\NormalTok{, }\FloatTok{1}\NormalTok{],}
\NormalTok{    xlabel}\OperatorTok{=}\StringTok{"antibiotic concentration"}\NormalTok{,}
\NormalTok{    ylabel}\OperatorTok{=}\StringTok{"optical density"}\NormalTok{,}
\NormalTok{    xscale}\OperatorTok{=}\NormalTok{log2}
\NormalTok{)}
\CommentTok{\# Plot data}
\FunctionTok{scatterlines!}\NormalTok{(ax, }\FloatTok{2.0} \OperatorTok{.\^{}}\NormalTok{ log2x, y)}

\NormalTok{fig}
\end{Highlighting}
\end{Shaded}

\includegraphics[width=7.29167in,height=6.25in]{iwasawa_bayesian_inference_files/figure-pdf/cell-9-output-1.png}

With the data simulated, let's perform inference on the data. For this,
we will use the NUTS sampler. To run multiple chains in parallel, we
will use the \texttt{MCMCThreads} option.

\begin{Shaded}
\begin{Highlighting}[]
\BuiltInTok{Random}\NormalTok{.}\FunctionTok{seed!}\NormalTok{(}\FloatTok{42}\NormalTok{)}
\CommentTok{\# Perform inference}
\NormalTok{model }\OperatorTok{=} \FunctionTok{logistic\_model}\NormalTok{(log2x, y)}

\CommentTok{\# Define number of steps}
\NormalTok{n\_burnin }\OperatorTok{=} \FloatTok{10\_000}
\NormalTok{n\_samples }\OperatorTok{=} \FloatTok{1\_000}

\CommentTok{\# Run sampler using}
\NormalTok{chain }\OperatorTok{=}\NormalTok{ Turing.}\FunctionTok{sample}\NormalTok{(}
\NormalTok{    model, Turing.}\FunctionTok{NUTS}\NormalTok{(), Turing.}\FunctionTok{MCMCThreads}\NormalTok{(), n\_burnin }\OperatorTok{+}\NormalTok{ n\_samples, }\FloatTok{4}
\NormalTok{)}
\end{Highlighting}
\end{Shaded}

\begin{verbatim}
Sampling (4 threads)   0%|                              |  ETA: N/A
┌ Info: Found initial step size
└   ϵ = 0.8
┌ Info: Found initial step size
└   ϵ = 0.8
┌ Info: Found initial step size
└   ϵ = 0.8
┌ Info: Found initial step size
└   ϵ = 0.4
Sampling (4 threads)  25%|███████▌                      |  ETA: 0:00:45
Sampling (4 threads)  50%|███████████████               |  ETA: 0:00:15
Sampling (4 threads)  75%|██████████████████████▌       |  ETA: 0:00:05
Sampling (4 threads) 100%|██████████████████████████████| Time: 0:00:14
Sampling (4 threads) 100%|██████████████████████████████| Time: 0:00:15
\end{verbatim}

\begin{verbatim}
Chains MCMC chain (11000×17×4 Array{Float64, 3}):

Iterations        = 1001:1:12000
Number of chains  = 4
Samples per chain = 11000
Wall duration     = 10.67 seconds
Compute duration  = 41.09 seconds
parameters        = log2ic50, a, b, c, σ²
internals         = lp, n_steps, is_accept, acceptance_rate, log_density, hamiltonian_energy, hamiltonian_energy_error, max_hamiltonian_energy_error, tree_depth, numerical_error, step_size, nom_step_size

Summary Statistics
  parameters      mean       std      mcse     ess_bulk     ess_tail      rhat ⋯
      Symbol   Float64   Float64   Float64      Float64      Float64   Float64 ⋯

    log2ic50    0.4723    0.0718    0.0005   23407.8672   14077.0993    1.0003 ⋯
           a    0.9668    0.0621    0.0005   16034.1809   16294.6985    1.0006 ⋯
           b    8.7902    4.6027    0.0449   13409.6206   13035.6999    1.0003 ⋯
           c    0.1192    0.0449    0.0004   15396.5130   11022.7335    1.0004 ⋯
          σ²    0.0111    0.0076    0.0001   12720.8496   16686.4143    1.0005 ⋯
                                                                1 column omitted

Quantiles
  parameters      2.5%     25.0%     50.0%     75.0%     97.5%
      Symbol   Float64   Float64   Float64   Float64   Float64

    log2ic50    0.3173    0.4352    0.4745    0.5141    0.6048
           a    0.8482    0.9279    0.9649    1.0042    1.0961
           b    3.0834    5.7843    7.9162   10.6434   20.2429
           c    0.0294    0.0902    0.1194    0.1475    0.2093
          σ²    0.0038    0.0065    0.0091    0.0132    0.0302
\end{verbatim}

Let's look at the posterior distribution of the parameters. For this, we
will use the \texttt{PairPlots.jl} package.

\begin{Shaded}
\begin{Highlighting}[]
\CommentTok{\# Plot corner plot for chains}
\NormalTok{PairPlots.}\FunctionTok{pairplot}\NormalTok{(}
\NormalTok{    chain[n\_burnin}\OperatorTok{+}\FloatTok{1}\OperatorTok{:}\KeywordTok{end}\NormalTok{, }\OperatorTok{:}\NormalTok{, }\OperatorTok{:}\NormalTok{],}
\NormalTok{    PairPlots.}\FunctionTok{Truth}\NormalTok{(}
\NormalTok{        (;}
\NormalTok{        log2ic50}\OperatorTok{=}\NormalTok{log2ic50\_true,}
\NormalTok{        a}\OperatorTok{=}\NormalTok{a\_true,}
\NormalTok{        b}\OperatorTok{=}\NormalTok{b\_true,}
\NormalTok{        c}\OperatorTok{=}\NormalTok{c\_true,}
\NormalTok{        σ²}\OperatorTok{=}\NormalTok{σ²\_true}
\NormalTok{    )}
\NormalTok{    )}
\NormalTok{)}
\end{Highlighting}
\end{Shaded}

\includegraphics[width=19.20833in,height=19.60417in]{iwasawa_bayesian_inference_files/figure-pdf/cell-11-output-1.png}

All parameters are well-constrained by the data, and we are able to
recover the ground truth values.

Let's plot the posterior predictive checks with the data to see how the
model performs.

\begin{Shaded}
\begin{Highlighting}[]
\BuiltInTok{Random}\NormalTok{.}\FunctionTok{seed!}\NormalTok{(}\FloatTok{42}\NormalTok{)}

\CommentTok{\# Initialize matrix to store samples}
\NormalTok{y\_samples }\OperatorTok{=} \FunctionTok{Array}\DataTypeTok{\{Float64\}}\NormalTok{(}\ConstantTok{undef}\NormalTok{, }\FunctionTok{length}\NormalTok{(log2x), n\_samples)}

\CommentTok{\# Loop through samples}
\ControlFlowTok{for}\NormalTok{ i }\KeywordTok{in} \FloatTok{1}\OperatorTok{:}\NormalTok{n\_samples}
    \CommentTok{\# Generate mean for sample}
\NormalTok{    y\_samples[}\OperatorTok{:}\NormalTok{, i] }\OperatorTok{=} \FunctionTok{logistic\_log2}\NormalTok{(}
\NormalTok{        log2x,}
\NormalTok{        chain[}\OperatorTok{:}\NormalTok{a][i],}
\NormalTok{        chain[}\OperatorTok{:}\NormalTok{b][i],}
\NormalTok{        chain[}\OperatorTok{:}\NormalTok{c][i],}
\NormalTok{        chain[}\OperatorTok{:}\NormalTok{log2ic50][i]}
\NormalTok{    )}
    \CommentTok{\# Add noise}
\NormalTok{    y\_samples[}\OperatorTok{:}\NormalTok{, i] }\OperatorTok{.+=} \FunctionTok{randn}\NormalTok{(}\FunctionTok{length}\NormalTok{(log2x)) }\OperatorTok{*} \FunctionTok{sqrt}\NormalTok{(chain[}\OperatorTok{:}\NormalTok{σ²][i])}
\ControlFlowTok{end} \CommentTok{\# for}

\CommentTok{\# Initialize figure}
\NormalTok{fig }\OperatorTok{=} \FunctionTok{Figure}\NormalTok{(size}\OperatorTok{=}\NormalTok{(}\FloatTok{350}\NormalTok{, }\FloatTok{300}\NormalTok{))}
\CommentTok{\# Add axis}
\NormalTok{ax }\OperatorTok{=} \FunctionTok{Axis}\NormalTok{(}
\NormalTok{    fig[}\FloatTok{1}\NormalTok{, }\FloatTok{1}\NormalTok{],}
\NormalTok{    xlabel}\OperatorTok{=}\StringTok{"antibiotic concentration"}\NormalTok{,}
\NormalTok{    ylabel}\OperatorTok{=}\StringTok{"optical density"}\NormalTok{,}
\NormalTok{    title}\OperatorTok{=}\StringTok{"Posterior Predictive Check"}\NormalTok{,}
\NormalTok{    xscale}\OperatorTok{=}\NormalTok{log2}
\NormalTok{)}

\CommentTok{\# Plot samples}
\ControlFlowTok{for}\NormalTok{ i }\KeywordTok{in} \FloatTok{1}\OperatorTok{:}\NormalTok{n\_samples}
    \FunctionTok{lines!}\NormalTok{(ax, }\FloatTok{2.0} \OperatorTok{.\^{}}\NormalTok{ log2x, y\_samples[}\OperatorTok{:}\NormalTok{, i], color}\OperatorTok{=}\NormalTok{(}\OperatorTok{:}\NormalTok{gray, }\FloatTok{0.05}\NormalTok{))}
\ControlFlowTok{end} \CommentTok{\# for}

\CommentTok{\# Plot data}
\FunctionTok{scatterlines!}\NormalTok{(ax, }\FloatTok{2.0} \OperatorTok{.\^{}}\NormalTok{ log2x, y)}

\NormalTok{fig}
\end{Highlighting}
\end{Shaded}

\includegraphics[width=7.29167in,height=6.25in]{iwasawa_bayesian_inference_files/figure-pdf/cell-12-output-1.png}

The fit looks excellent. Let's now perform inference on the Iwasawa
data. We will use one example dataset. Moreover, we will set informative
priors for the parameters.

\begin{Shaded}
\begin{Highlighting}[]
\BuiltInTok{Random}\NormalTok{.}\FunctionTok{seed!}\NormalTok{(}\FloatTok{42}\NormalTok{)}
\CommentTok{\# Group data by antibiotic, environment, and day}
\NormalTok{df\_group }\OperatorTok{=}\NormalTok{ DF.}\FunctionTok{groupby}\NormalTok{(}
\NormalTok{    df[(.!df.blank)}\OperatorTok{.\&}\NormalTok{(df.concentration\_ugmL}\OperatorTok{.\textgreater{}}\FloatTok{0}\NormalTok{), }\OperatorTok{:}\NormalTok{],}
\NormalTok{    [}\OperatorTok{:}\NormalTok{antibiotic, }\OperatorTok{:}\NormalTok{env, }\OperatorTok{:}\NormalTok{day]}
\NormalTok{)}

\CommentTok{\# Extract data}
\NormalTok{data }\OperatorTok{=}\NormalTok{ df\_group[}\FloatTok{2}\NormalTok{]}

\CommentTok{\# Define prior parameters}
\NormalTok{prior\_params }\OperatorTok{=}\NormalTok{ (}
\NormalTok{    a}\OperatorTok{=}\NormalTok{(}\FunctionTok{log}\NormalTok{(}\FloatTok{0.1}\NormalTok{), }\FloatTok{0.1}\NormalTok{),}
\NormalTok{    b}\OperatorTok{=}\NormalTok{(}\FloatTok{0}\NormalTok{, }\FloatTok{1}\NormalTok{),}
\NormalTok{    c}\OperatorTok{=}\NormalTok{(}\FloatTok{0}\NormalTok{, }\FloatTok{1}\NormalTok{),}
\NormalTok{    σ²}\OperatorTok{=}\NormalTok{(}\FloatTok{0}\NormalTok{, }\FloatTok{0.1}\NormalTok{)}
\NormalTok{)}
\CommentTok{\# Perform inference}
\NormalTok{model }\OperatorTok{=} \FunctionTok{logistic\_model}\NormalTok{(}
    \FunctionTok{log2}\NormalTok{.(data.concentration\_ugmL),}
\NormalTok{    data.OD,}
\NormalTok{    prior\_params}
\NormalTok{)}

\CommentTok{\# Define number of steps}
\NormalTok{n\_burnin }\OperatorTok{=} \FloatTok{10\_000}
\NormalTok{n\_samples }\OperatorTok{=} \FloatTok{1\_000}

\NormalTok{chain }\OperatorTok{=}\NormalTok{ Turing.}\FunctionTok{sample}\NormalTok{(}
\NormalTok{    model, Turing.}\FunctionTok{NUTS}\NormalTok{(), Turing.}\FunctionTok{MCMCThreads}\NormalTok{(), n\_burnin }\OperatorTok{+}\NormalTok{ n\_samples, }\FloatTok{4}
\NormalTok{)}
\end{Highlighting}
\end{Shaded}

\begin{verbatim}
Sampling (4 threads)   0%|                              |  ETA: N/A
┌ Info: Found initial step size
└   ϵ = 0.2
┌ Info: Found initial step size
└   ϵ = 0.05
┌ Info: Found initial step size
└   ϵ = 0.2
┌ Info: Found initial step size
└   ϵ = 0.05
Sampling (4 threads)  25%|███████▌                      |  ETA: 0:00:17
Sampling (4 threads)  50%|███████████████               |  ETA: 0:00:06
Sampling (4 threads)  75%|██████████████████████▌       |  ETA: 0:00:02
Sampling (4 threads) 100%|██████████████████████████████| Time: 0:00:05
Sampling (4 threads) 100%|██████████████████████████████| Time: 0:00:05
\end{verbatim}

\begin{verbatim}
Chains MCMC chain (11000×17×4 Array{Float64, 3}):

Iterations        = 1001:1:12000
Number of chains  = 4
Samples per chain = 11000
Wall duration     = 5.05 seconds
Compute duration  = 19.36 seconds
parameters        = log2ic50, a, b, c, σ²
internals         = lp, n_steps, is_accept, acceptance_rate, log_density, hamiltonian_energy, hamiltonian_energy_error, max_hamiltonian_energy_error, tree_depth, numerical_error, step_size, nom_step_size

Summary Statistics
  parameters      mean       std      mcse     ess_bulk     ess_tail      rhat ⋯
      Symbol   Float64   Float64   Float64      Float64      Float64   Float64 ⋯

    log2ic50    3.1585    0.0611    0.0004   32325.5670   25247.0891    1.0001 ⋯
           a    0.1013    0.0074    0.0000   28541.0328   27631.5354    1.0001 ⋯
           b   15.8931   10.4326    0.0655   31581.5518   26387.0900    1.0000 ⋯
           c    0.0873    0.0073    0.0000   27591.1350   24819.1854    1.0002 ⋯
          σ²    0.0015    0.0002    0.0000   37231.0264   30213.0584    1.0000 ⋯
                                                                1 column omitted

Quantiles
  parameters      2.5%     25.0%     50.0%     75.0%     97.5%
      Symbol   Float64   Float64   Float64   Float64   Float64

    log2ic50    3.0427    3.1192    3.1589    3.1963    3.2817
           a    0.0872    0.0962    0.1012    0.1063    0.1162
           b    4.6324    9.2475   13.4282   19.5750   41.4435
           c    0.0726    0.0824    0.0874    0.0923    0.1013
          σ²    0.0011    0.0013    0.0014    0.0016    0.0020
\end{verbatim}

Let's look at the corner plot for the chains.

\begin{Shaded}
\begin{Highlighting}[]
\CommentTok{\# Plot corner plot for chains}
\NormalTok{PairPlots.}\FunctionTok{pairplot}\NormalTok{(chain[n\_burnin}\OperatorTok{+}\FloatTok{1}\OperatorTok{:}\KeywordTok{end}\NormalTok{, }\OperatorTok{:}\NormalTok{, }\OperatorTok{:}\NormalTok{])}
\end{Highlighting}
\end{Shaded}

\includegraphics[width=19.4375in,height=19.83333in]{iwasawa_bayesian_inference_files/figure-pdf/cell-14-output-1.png}

This looks good. Let's now plot the posterior predictive check.

\begin{Shaded}
\begin{Highlighting}[]
\BuiltInTok{Random}\NormalTok{.}\FunctionTok{seed!}\NormalTok{(}\FloatTok{42}\NormalTok{)}

\CommentTok{\# Define unique concentrations}
\NormalTok{unique\_concentrations }\OperatorTok{=} \FunctionTok{sort}\NormalTok{(}\FunctionTok{unique}\NormalTok{(data.concentration\_ugmL))}

\CommentTok{\# Initialize matrix to store samples}
\NormalTok{y\_samples }\OperatorTok{=} \FunctionTok{Array}\DataTypeTok{\{Float64\}}\NormalTok{(}
    \ConstantTok{undef}\NormalTok{, }\FunctionTok{length}\NormalTok{(unique\_concentrations), n\_samples}
\NormalTok{)}

\CommentTok{\# Loop through samples}
\ControlFlowTok{for}\NormalTok{ i }\KeywordTok{in} \FloatTok{1}\OperatorTok{:}\NormalTok{n\_samples}
    \CommentTok{\# Generate mean for sample}
\NormalTok{    y\_samples[}\OperatorTok{:}\NormalTok{, i] }\OperatorTok{=} \FunctionTok{logistic\_log2}\NormalTok{(}
        \FunctionTok{log2}\NormalTok{.(unique\_concentrations),}
\NormalTok{        chain[}\OperatorTok{:}\NormalTok{a][i],}
\NormalTok{        chain[}\OperatorTok{:}\NormalTok{b][i],}
\NormalTok{        chain[}\OperatorTok{:}\NormalTok{c][i],}
\NormalTok{        chain[}\OperatorTok{:}\NormalTok{log2ic50][i]}
\NormalTok{    )}
    \CommentTok{\# Add noise}
\NormalTok{    y\_samples[}\OperatorTok{:}\NormalTok{, i] }\OperatorTok{.+=} \FunctionTok{randn}\NormalTok{(}\FunctionTok{length}\NormalTok{(unique\_concentrations)) }\OperatorTok{*} \FunctionTok{√}\NormalTok{(chain[}\OperatorTok{:}\NormalTok{σ²][i])}
    \FunctionTok{√}\NormalTok{(chain[}\OperatorTok{:}\NormalTok{σ²][i])}
\ControlFlowTok{end} \CommentTok{\# for}

\CommentTok{\# Initialize figure}
\NormalTok{fig }\OperatorTok{=} \FunctionTok{Figure}\NormalTok{(size}\OperatorTok{=}\NormalTok{(}\FloatTok{350}\NormalTok{, }\FloatTok{300}\NormalTok{))}
\CommentTok{\# Add axis}
\NormalTok{ax }\OperatorTok{=} \FunctionTok{Axis}\NormalTok{(}
\NormalTok{    fig[}\FloatTok{1}\NormalTok{, }\FloatTok{1}\NormalTok{],}
\NormalTok{    xlabel}\OperatorTok{=}\StringTok{"antibiotic concentration"}\NormalTok{,}
\NormalTok{    ylabel}\OperatorTok{=}\StringTok{"optical density"}\NormalTok{,}
\NormalTok{    title}\OperatorTok{=}\StringTok{"Posterior Predictive Check"}\NormalTok{,}
\NormalTok{    xscale}\OperatorTok{=}\NormalTok{log2}
\NormalTok{)}

\CommentTok{\# Plot samples}
\ControlFlowTok{for}\NormalTok{ i }\KeywordTok{in} \FloatTok{1}\OperatorTok{:}\NormalTok{n\_samples}
    \FunctionTok{lines!}\NormalTok{(}
\NormalTok{        ax,}
\NormalTok{        unique\_concentrations,}
\NormalTok{        y\_samples[}\OperatorTok{:}\NormalTok{, i],}
\NormalTok{        color}\OperatorTok{=}\NormalTok{(ColorSchemes.Paired\_12[}\FloatTok{1}\NormalTok{], }\FloatTok{0.5}\NormalTok{)}
\NormalTok{    )}
\ControlFlowTok{end} \CommentTok{\# for}

\CommentTok{\# Plot data}
\FunctionTok{scatter!}\NormalTok{(ax, data.concentration\_ugmL, data.OD)}

\NormalTok{fig}
\end{Highlighting}
\end{Shaded}

\includegraphics[width=7.29167in,height=6.25in]{iwasawa_bayesian_inference_files/figure-pdf/cell-15-output-1.png}

The presence of the outlier measurements expands the posterior
predictive checks uncertainty.

Let's look into how to detect outliers in the data.

\subsection{Residual-based outlier
detection}\label{residual-based-outlier-detection}

Our naive approach to detect outliers is to fit the logistic model
deterministically and then identify points that are more than 3 standard
deviations from the mean. This is known as ``residual-based outlier
detection''.

\begin{Shaded}
\begin{Highlighting}[]
\KeywordTok{function} \FunctionTok{logistic\_log2}\NormalTok{(log2x, params)}
    \ControlFlowTok{return} \FunctionTok{logistic\_log2}\NormalTok{(log2x, params}\OperatorTok{...}\NormalTok{)}
\KeywordTok{end}

\StringTok{"""}
\StringTok{    fit\_logistic\_and\_detect\_outliers(log2x, y; threshold=3)}

\StringTok{Fit a logistic model to the given data and detect outliers based on residuals.}

\StringTok{This function performs the following steps:}
\StringTok{1. Fits a logistic model to the log2{-}transformed x{-}values and y{-}values.}
\StringTok{2. Calculates residuals between the fitted model and actual y{-}values.}
\StringTok{3. Identifies outliers as points with residuals exceeding a specified threshold.}

\StringTok{\# Arguments}
\StringTok{{-} \textasciigrave{}log2x\textasciigrave{}: Array of log2{-}transformed x{-}values (typically concentrations).}
\StringTok{{-} \textasciigrave{}y\textasciigrave{}: Array of y{-}values (typically optical density measurements).}
\StringTok{{-} \textasciigrave{}threshold\textasciigrave{}: Number of standard deviations beyond which a point is considered an outlier. Default is 3.}

\StringTok{\# Returns}
\StringTok{{-} A boolean array indicating which points are outliers (true for outliers).}

\StringTok{\# Notes}
\StringTok{{-} The function uses a logistic model of the form: }
\StringTok{    y = a / (1 + exp(b * (log2x {-} log2ic50))) + c}
\StringTok{{-} Initial parameter guesses are made based on the input data.}
\StringTok{{-} The LsqFit package is used for curve fitting.}
\StringTok{{-} Outliers are determined by comparing the absolute residuals to the threshold * standard deviation of residuals.}
\StringTok{"""}
\KeywordTok{function} \FunctionTok{fit\_logistic\_and\_detect\_outliers}\NormalTok{(log2x, y; threshold}\OperatorTok{=}\FloatTok{3}\NormalTok{)}
    \CommentTok{\# Initial parameter guess}
\NormalTok{    p0 }\OperatorTok{=}\NormalTok{ [}\FloatTok{0.1}\NormalTok{, }\FloatTok{1.0}\NormalTok{, }\FunctionTok{maximum}\NormalTok{(y) }\OperatorTok{{-}} \FunctionTok{minimum}\NormalTok{(y), StatsBase.}\FunctionTok{median}\NormalTok{(log2x)]}

    \CommentTok{\# Fit the logistic model}
\NormalTok{    fit }\OperatorTok{=}\NormalTok{ LsqFit.}\FunctionTok{curve\_fit}\NormalTok{(logistic\_log2, log2x, y, p0)}

    \CommentTok{\# Calculate residuals}
\NormalTok{    residuals }\OperatorTok{=}\NormalTok{ y }\OperatorTok{{-}} \FunctionTok{logistic\_log2}\NormalTok{(log2x, fit.param)}

    \CommentTok{\# Calculate standard deviation of residuals}
\NormalTok{    σ }\OperatorTok{=}\NormalTok{ StatsBase.}\FunctionTok{std}\NormalTok{(residuals)}

    \CommentTok{\# Identify outliers}
\NormalTok{    outliers\_idx }\OperatorTok{=} \FunctionTok{abs}\NormalTok{.(residuals) }\OperatorTok{.\textgreater{}}\NormalTok{ threshold }\OperatorTok{*}\NormalTok{ σ}

    \CommentTok{\# Return outlier indices}
    \ControlFlowTok{return}\NormalTok{ outliers\_idx}
\KeywordTok{end}
\end{Highlighting}
\end{Shaded}

Let's test this function by trying to remove the outliers from the data
we used in the previous section.

\begin{Shaded}
\begin{Highlighting}[]
\CommentTok{\# Locate outliers}
\NormalTok{outliers\_idx }\OperatorTok{=} \FunctionTok{fit\_logistic\_and\_detect\_outliers}\NormalTok{(}
    \FunctionTok{log2}\NormalTok{.(data.concentration\_ugmL), data.OD, threshold}\OperatorTok{=}\FloatTok{2}
\NormalTok{)}

\CommentTok{\# Plot the results}
\NormalTok{fig }\OperatorTok{=} \FunctionTok{Figure}\NormalTok{(size}\OperatorTok{=}\NormalTok{(}\FloatTok{450}\NormalTok{, }\FloatTok{300}\NormalTok{))}

\NormalTok{ax }\OperatorTok{=} \FunctionTok{Axis}\NormalTok{(}
\NormalTok{    fig[}\FloatTok{1}\NormalTok{, }\FloatTok{1}\NormalTok{],}
\NormalTok{    xlabel}\OperatorTok{=}\StringTok{"antibiotic concentration"}\NormalTok{,}
\NormalTok{    ylabel}\OperatorTok{=}\StringTok{"optical density"}\NormalTok{,}
    \CommentTok{\# xscale=log2}
\NormalTok{)}

\CommentTok{\# Plot the original data}
\FunctionTok{scatter!}\NormalTok{(}
\NormalTok{    ax,}
\NormalTok{    data.concentration\_ugmL,}
\NormalTok{    data.OD,}
\NormalTok{    color}\OperatorTok{=}\NormalTok{ColorSchemes.Paired\_12[}\FloatTok{1}\NormalTok{],}
\NormalTok{    label}\OperatorTok{=}\StringTok{"data"}
\NormalTok{)}

\CommentTok{\# Plot the cleaned data}
\FunctionTok{scatter!}\NormalTok{(}
\NormalTok{    ax,}
\NormalTok{    data.concentration\_ugmL[.!outliers\_idx] }\OperatorTok{.+} \FloatTok{0.5}\NormalTok{,}
\NormalTok{    data.OD[.!outliers\_idx], color}\OperatorTok{=}\NormalTok{ColorSchemes.Paired\_12[}\FloatTok{2}\NormalTok{],}
\NormalTok{    label}\OperatorTok{=}\StringTok{"cleaned data"}
\NormalTok{)}

\CommentTok{\# Add legend}
\FunctionTok{Legend}\NormalTok{(fig[}\FloatTok{1}\NormalTok{, }\FloatTok{2}\NormalTok{], ax)}

\NormalTok{fig}
\end{Highlighting}
\end{Shaded}

\includegraphics[width=9.375in,height=6.25in]{iwasawa_bayesian_inference_files/figure-pdf/cell-17-output-1.png}

The detection of outliers worked really well. Let's now perform
inference on the cleaned data.

\begin{Shaded}
\begin{Highlighting}[]
\BuiltInTok{Random}\NormalTok{.}\FunctionTok{seed!}\NormalTok{(}\FloatTok{42}\NormalTok{)}
\CommentTok{\# Group data by antibiotic, environment, and day}
\NormalTok{df\_group }\OperatorTok{=}\NormalTok{ DF.}\FunctionTok{groupby}\NormalTok{(}
\NormalTok{    df[(.!df.blank)}\OperatorTok{.\&}\NormalTok{(df.concentration\_ugmL}\OperatorTok{.\textgreater{}}\FloatTok{0}\NormalTok{), }\OperatorTok{:}\NormalTok{],}
\NormalTok{    [}\OperatorTok{:}\NormalTok{antibiotic, }\OperatorTok{:}\NormalTok{env, }\OperatorTok{:}\NormalTok{day]}
\NormalTok{)}

\CommentTok{\# Extract data}
\NormalTok{data }\OperatorTok{=}\NormalTok{ df\_group[}\FloatTok{2}\NormalTok{]}

\CommentTok{\# Find outliers}
\NormalTok{outliers\_idx }\OperatorTok{=} \FunctionTok{fit\_logistic\_and\_detect\_outliers}\NormalTok{(}
    \FunctionTok{log2}\NormalTok{.(data.concentration\_ugmL), data.OD, threshold}\OperatorTok{=}\FloatTok{2}
\NormalTok{)}
\CommentTok{\# Remove outliers}
\NormalTok{data\_clean }\OperatorTok{=}\NormalTok{ data[.!outliers\_idx, }\OperatorTok{:}\NormalTok{]}

\CommentTok{\# Define prior parameters}
\NormalTok{prior\_params }\OperatorTok{=}\NormalTok{ (}
\NormalTok{    a}\OperatorTok{=}\NormalTok{(}\FunctionTok{log}\NormalTok{(}\FloatTok{0.1}\NormalTok{), }\FloatTok{0.1}\NormalTok{),}
\NormalTok{    b}\OperatorTok{=}\NormalTok{(}\FloatTok{0}\NormalTok{, }\FloatTok{1}\NormalTok{),}
\NormalTok{    c}\OperatorTok{=}\NormalTok{(}\FloatTok{0}\NormalTok{, }\FloatTok{1}\NormalTok{),}
\NormalTok{    σ²}\OperatorTok{=}\NormalTok{(}\FloatTok{0}\NormalTok{, }\FloatTok{0.1}\NormalTok{)}
\NormalTok{)}

\CommentTok{\# Perform inference}
\NormalTok{model }\OperatorTok{=} \FunctionTok{logistic\_model}\NormalTok{(}
    \FunctionTok{log2}\NormalTok{.(data\_clean.concentration\_ugmL),}
\NormalTok{    data\_clean.OD,}
\NormalTok{    prior\_params}
\NormalTok{)}

\CommentTok{\# Define number of steps}
\NormalTok{n\_burnin }\OperatorTok{=} \FloatTok{10\_000}
\NormalTok{n\_samples }\OperatorTok{=} \FloatTok{1\_000}

\NormalTok{chain }\OperatorTok{=}\NormalTok{ Turing.}\FunctionTok{sample}\NormalTok{(}
\NormalTok{    model, Turing.}\FunctionTok{NUTS}\NormalTok{(), Turing.}\FunctionTok{MCMCThreads}\NormalTok{(), n\_burnin }\OperatorTok{+}\NormalTok{ n\_samples, }\FloatTok{4}
\NormalTok{)}
\end{Highlighting}
\end{Shaded}

\begin{verbatim}
Sampling (4 threads)   0%|                              |  ETA: N/A
┌ Info: Found initial step size
└   ϵ = 0.05
┌ Info: Found initial step size
└   ϵ = 0.2
┌ Info: Found initial step size
└   ϵ = 0.2
┌ Info: Found initial step size
└   ϵ = 0.05
Sampling (4 threads)  25%|███████▌                      |  ETA: 0:00:18
Sampling (4 threads)  50%|███████████████               |  ETA: 0:00:06
Sampling (4 threads)  75%|██████████████████████▌       |  ETA: 0:00:02
Sampling (4 threads) 100%|██████████████████████████████| Time: 0:00:06
Sampling (4 threads) 100%|██████████████████████████████| Time: 0:00:06
\end{verbatim}

\begin{verbatim}
Chains MCMC chain (11000×17×4 Array{Float64, 3}):

Iterations        = 1001:1:12000
Number of chains  = 4
Samples per chain = 11000
Wall duration     = 5.38 seconds
Compute duration  = 20.29 seconds
parameters        = log2ic50, a, b, c, σ²
internals         = lp, n_steps, is_accept, acceptance_rate, log_density, hamiltonian_energy, hamiltonian_energy_error, max_hamiltonian_energy_error, tree_depth, numerical_error, step_size, nom_step_size

Summary Statistics
  parameters      mean       std      mcse     ess_bulk     ess_tail      rhat ⋯
      Symbol   Float64   Float64   Float64      Float64      Float64   Float64 ⋯

    log2ic50    3.1001    0.0297    0.0002   25150.3685   27935.5938    1.0002 ⋯
           a    0.1163    0.0038    0.0000   16860.1510   22570.9089    1.0001 ⋯
           b    6.7350    1.0786    0.0069   26759.0743   25218.0167    1.0001 ⋯
           c    0.0712    0.0034    0.0000   16323.0552   21352.9296    1.0001 ⋯
          σ²    0.0001    0.0000    0.0000   32838.7703   28866.3125    1.0001 ⋯
                                                                1 column omitted

Quantiles
  parameters      2.5%     25.0%     50.0%     75.0%     97.5%
      Symbol   Float64   Float64   Float64   Float64   Float64

    log2ic50    3.0410    3.0804    3.1004    3.1202    3.1580
           a    0.1089    0.1138    0.1163    0.1188    0.1238
           b    5.0241    5.9818    6.5979    7.3261    9.2492
           c    0.0643    0.0689    0.0712    0.0735    0.0778
          σ²    0.0001    0.0001    0.0001    0.0001    0.0002
\end{verbatim}

Let's look again at the corner plot for the chains.

\begin{Shaded}
\begin{Highlighting}[]
\CommentTok{\# Plot corner plot for chains}
\NormalTok{PairPlots.}\FunctionTok{pairplot}\NormalTok{(chain[n\_burnin}\OperatorTok{+}\FloatTok{1}\OperatorTok{:}\KeywordTok{end}\NormalTok{, }\OperatorTok{:}\NormalTok{, }\OperatorTok{:}\NormalTok{])}
\end{Highlighting}
\end{Shaded}

\includegraphics[width=19.5625in,height=20.02083in]{iwasawa_bayesian_inference_files/figure-pdf/cell-19-output-1.png}

Compared to the previous example where we didn't remove the outliers,
the posterior distributions are more concentrated, especially for the
\(b\) parameter.

Let's now plot the posterior predictive check.

\begin{Shaded}
\begin{Highlighting}[]
\BuiltInTok{Random}\NormalTok{.}\FunctionTok{seed!}\NormalTok{(}\FloatTok{42}\NormalTok{)}

\CommentTok{\# Define unique concentrations}
\NormalTok{unique\_concentrations }\OperatorTok{=} \FunctionTok{sort}\NormalTok{(}\FunctionTok{unique}\NormalTok{(data\_clean.concentration\_ugmL))}

\CommentTok{\# Initialize matrix to store samples}
\NormalTok{y\_samples }\OperatorTok{=} \FunctionTok{Array}\DataTypeTok{\{Float64\}}\NormalTok{(}
    \ConstantTok{undef}\NormalTok{, }\FunctionTok{length}\NormalTok{(unique\_concentrations), n\_samples}
\NormalTok{)}

\CommentTok{\# Loop through samples}
\ControlFlowTok{for}\NormalTok{ i }\KeywordTok{in} \FloatTok{1}\OperatorTok{:}\NormalTok{n\_samples}
    \CommentTok{\# Generate mean for sample}
\NormalTok{    y\_samples[}\OperatorTok{:}\NormalTok{, i] }\OperatorTok{=} \FunctionTok{logistic\_log2}\NormalTok{(}
        \FunctionTok{log2}\NormalTok{.(unique\_concentrations),}
\NormalTok{        chain[}\OperatorTok{:}\NormalTok{a][i],}
\NormalTok{        chain[}\OperatorTok{:}\NormalTok{b][i],}
\NormalTok{        chain[}\OperatorTok{:}\NormalTok{c][i],}
\NormalTok{        chain[}\OperatorTok{:}\NormalTok{log2ic50][i]}
\NormalTok{    )}
    \CommentTok{\# Add noise}
\NormalTok{    y\_samples[}\OperatorTok{:}\NormalTok{, i] }\OperatorTok{.+=} \FunctionTok{randn}\NormalTok{(}\FunctionTok{length}\NormalTok{(unique\_concentrations)) }\OperatorTok{*} \FunctionTok{√}\NormalTok{(chain[}\OperatorTok{:}\NormalTok{σ²][i])}
    \FunctionTok{√}\NormalTok{(chain[}\OperatorTok{:}\NormalTok{σ²][i])}
\ControlFlowTok{end} \CommentTok{\# for}

\CommentTok{\# Initialize figure}
\NormalTok{fig }\OperatorTok{=} \FunctionTok{Figure}\NormalTok{(size}\OperatorTok{=}\NormalTok{(}\FloatTok{350}\NormalTok{, }\FloatTok{300}\NormalTok{))}
\CommentTok{\# Add axis}
\NormalTok{ax }\OperatorTok{=} \FunctionTok{Axis}\NormalTok{(}
\NormalTok{    fig[}\FloatTok{1}\NormalTok{, }\FloatTok{1}\NormalTok{],}
\NormalTok{    xlabel}\OperatorTok{=}\StringTok{"antibiotic concentration"}\NormalTok{,}
\NormalTok{    ylabel}\OperatorTok{=}\StringTok{"optical density"}\NormalTok{,}
\NormalTok{    title}\OperatorTok{=}\StringTok{"Posterior Predictive Check"}\NormalTok{,}
\NormalTok{    xscale}\OperatorTok{=}\NormalTok{log2}
\NormalTok{)}

\CommentTok{\# Plot samples}
\ControlFlowTok{for}\NormalTok{ i }\KeywordTok{in} \FloatTok{1}\OperatorTok{:}\NormalTok{n\_samples}
    \FunctionTok{lines!}\NormalTok{(}
\NormalTok{        ax,}
\NormalTok{        unique\_concentrations,}
\NormalTok{        y\_samples[}\OperatorTok{:}\NormalTok{, i],}
\NormalTok{        color}\OperatorTok{=}\NormalTok{(ColorSchemes.Paired\_12[}\FloatTok{1}\NormalTok{], }\FloatTok{0.5}\NormalTok{)}
\NormalTok{    )}
\ControlFlowTok{end} \CommentTok{\# for}

\CommentTok{\# Plot data}
\FunctionTok{scatter!}\NormalTok{(ax, data\_clean.concentration\_ugmL, data\_clean.OD)}

\NormalTok{fig}
\end{Highlighting}
\end{Shaded}

\includegraphics[width=7.29167in,height=6.25in]{iwasawa_bayesian_inference_files/figure-pdf/cell-20-output-1.png}

This look much better. Removing the outliers improved the fit
significantly.

\subsection{Alternative
parameterization}\label{alternative-parameterization}

In the data, there are several concentrations at zero, which is not
compatible with the log-scale model. However, we can rewrite
Equation~\ref{eq-iwasawa-model} without using the exponent as

\begin{equation}\phantomsection\label{eq-iwasawa-model-no-exponent}{
f(x) = \frac{a}
{1+ \left(\frac{x}{\mathrm{IC}_{50}}\right)^b} + c
}\end{equation}

Let's define a new model using this parameterization.

\begin{Shaded}
\begin{Highlighting}[]
\PreprocessorTok{@doc} \StringTok{raw"""}
\StringTok{    logistic\_alt(x, a, b, c, ic50)}

\StringTok{Compute the logistic function used to model the relationship between antibiotic}
\StringTok{concentration and bacterial growth.}

\StringTok{This function implements the following equation:}

\StringTok{f(x) = a / (1 + (x / IC₅₀) \^{} b) + c}

\StringTok{\# Arguments}
\StringTok{{-} \textasciigrave{}x\textasciigrave{}: Antibiotic concentration (input variable)}
\StringTok{{-} \textasciigrave{}a\textasciigrave{}: Maximum effect parameter (difference between upper and lower asymptotes)}
\StringTok{{-} \textasciigrave{}b\textasciigrave{}: Slope parameter (steepness of the curve)}
\StringTok{{-} \textasciigrave{}c\textasciigrave{}: Minimum effect parameter (lower asymptote)}
\StringTok{{-} \textasciigrave{}ic50\textasciigrave{}: IC₅₀ parameter (concentration at which the effect is halfway between}
\StringTok{  the minimum and maximum)}

\StringTok{\# Returns}
\StringTok{The computed effect (e.g., optical density) for the given antibiotic}
\StringTok{concentration and parameters.}

\StringTok{Note: This function is vectorized and can handle array inputs for \textasciigrave{}x\textasciigrave{}.}
\StringTok{"""}
\KeywordTok{function} \FunctionTok{logistic\_alt}\NormalTok{(x, a, b, c, ic50)}
    \ControlFlowTok{return}\NormalTok{ @. a }\OperatorTok{/}\NormalTok{ (}\FloatTok{1.0} \OperatorTok{+}\NormalTok{ (x }\OperatorTok{/}\NormalTok{ ic50)}\OperatorTok{\^{}}\NormalTok{b) }\OperatorTok{+}\NormalTok{ c}
\KeywordTok{end}

\KeywordTok{function} \FunctionTok{logistic\_alt}\NormalTok{(x, params)}
    \ControlFlowTok{return} \FunctionTok{logistic\_alt}\NormalTok{(x, params}\OperatorTok{...}\NormalTok{)}
\KeywordTok{end}
\end{Highlighting}
\end{Shaded}

We will use the previous function to detect outliers in the data. Now,
we re-define the Bayesian model using this parameterization.

\begin{Shaded}
\begin{Highlighting}[]
\NormalTok{Turing.}\PreprocessorTok{@model} \KeywordTok{function} \FunctionTok{logistic\_alt\_model}\NormalTok{(}
\NormalTok{    x, y, prior\_params}\OperatorTok{::}\DataTypeTok{NamedTuple}\NormalTok{=}\FunctionTok{NamedTuple}\NormalTok{()}
\NormalTok{)}
    \CommentTok{\# Define default prior parameters}
\NormalTok{    default\_params }\OperatorTok{=}\NormalTok{ (}
\NormalTok{        ic50}\OperatorTok{=}\NormalTok{(}\FloatTok{0}\NormalTok{, }\FloatTok{1}\NormalTok{),}
\NormalTok{        a}\OperatorTok{=}\NormalTok{(}\FloatTok{0}\NormalTok{, }\FloatTok{1}\NormalTok{),}
\NormalTok{        b}\OperatorTok{=}\NormalTok{(}\FloatTok{0}\NormalTok{, }\FloatTok{1}\NormalTok{),}
\NormalTok{        c}\OperatorTok{=}\NormalTok{(}\FloatTok{0}\NormalTok{, }\FloatTok{1}\NormalTok{),}
\NormalTok{        σ²}\OperatorTok{=}\NormalTok{(}\FloatTok{0}\NormalTok{, }\FloatTok{1}\NormalTok{)}
\NormalTok{    )}

    \CommentTok{\# Merge default parameters with provided parameters}
\NormalTok{    params }\OperatorTok{=} \FunctionTok{merge}\NormalTok{(default\_params, prior\_params)}

    \CommentTok{\# Define priors}
\NormalTok{    ic50 }\OperatorTok{\textasciitilde{}}\NormalTok{ Turing.}\FunctionTok{LogNormal}\NormalTok{(params.ic50}\OperatorTok{...}\NormalTok{)}
\NormalTok{    a }\OperatorTok{\textasciitilde{}}\NormalTok{ Turing.}\FunctionTok{LogNormal}\NormalTok{(params.a}\OperatorTok{...}\NormalTok{)}
\NormalTok{    b }\OperatorTok{\textasciitilde{}}\NormalTok{ Turing.}\FunctionTok{LogNormal}\NormalTok{(params.b}\OperatorTok{...}\NormalTok{)}
\NormalTok{    c }\OperatorTok{\textasciitilde{}}\NormalTok{ Turing.}\FunctionTok{truncated}\NormalTok{(Turing.}\FunctionTok{Normal}\NormalTok{(params.c}\OperatorTok{...}\NormalTok{), }\FloatTok{0}\NormalTok{, }\ConstantTok{Inf}\NormalTok{)}
\NormalTok{    σ² }\OperatorTok{\textasciitilde{}}\NormalTok{ Turing.}\FunctionTok{truncated}\NormalTok{(Turing.}\FunctionTok{Normal}\NormalTok{(params.σ²}\OperatorTok{...}\NormalTok{), }\FloatTok{0}\NormalTok{, }\ConstantTok{Inf}\NormalTok{)}

    \CommentTok{\# Define likelihood}
\NormalTok{    y }\OperatorTok{\textasciitilde{}}\NormalTok{ Turing.}\FunctionTok{MvNormal}\NormalTok{(}
        \FunctionTok{logistic\_alt}\NormalTok{(x, a, b, c, ic50),}
        \BuiltInTok{LinearAlgebra}\NormalTok{.}\FunctionTok{Diagonal}\NormalTok{(}\FunctionTok{fill}\NormalTok{(σ², }\FunctionTok{length}\NormalTok{(y)))}
\NormalTok{    )}
\KeywordTok{end}
\end{Highlighting}
\end{Shaded}

Let's now perform inference on synthetic data using this new
parameterization.

\begin{Shaded}
\begin{Highlighting}[]
\BuiltInTok{Random}\NormalTok{.}\FunctionTok{seed!}\NormalTok{(}\FloatTok{42}\NormalTok{)}
\CommentTok{\# Define ground truth parameters}
\NormalTok{ic50\_true }\OperatorTok{=} \FloatTok{2}\OperatorTok{\^{}}\FloatTok{0.5}
\NormalTok{a\_true }\OperatorTok{=} \FloatTok{1.0}
\NormalTok{b\_true }\OperatorTok{=} \FloatTok{10.0}
\NormalTok{σ²\_true }\OperatorTok{=} \FloatTok{0.01}
\NormalTok{c\_true }\OperatorTok{=} \FloatTok{0.1}

\CommentTok{\# Simulate data}
\NormalTok{x }\OperatorTok{=} \FloatTok{2} \OperatorTok{.\^{}} \FunctionTok{LinRange}\NormalTok{(}\OperatorTok{{-}}\FloatTok{2.5}\NormalTok{, }\FloatTok{2.5}\NormalTok{, }\FloatTok{15}\NormalTok{)}
\CommentTok{\# Define mean of data}
\NormalTok{ŷ }\OperatorTok{=} \FunctionTok{logistic\_alt}\NormalTok{(x, a\_true, b\_true, c\_true, ic50\_true)}
\CommentTok{\# Add noise}
\NormalTok{y }\OperatorTok{=}\NormalTok{ ŷ }\OperatorTok{+} \FunctionTok{randn}\NormalTok{(}\FunctionTok{length}\NormalTok{(x)) }\OperatorTok{*} \FunctionTok{sqrt}\NormalTok{(σ²\_true)}

\CommentTok{\# Initialize figure}
\NormalTok{fig }\OperatorTok{=} \FunctionTok{Figure}\NormalTok{(size}\OperatorTok{=}\NormalTok{(}\FloatTok{350}\NormalTok{, }\FloatTok{300}\NormalTok{))}
\CommentTok{\# Add axis}
\NormalTok{ax }\OperatorTok{=} \FunctionTok{Axis}\NormalTok{(}
\NormalTok{    fig[}\FloatTok{1}\NormalTok{, }\FloatTok{1}\NormalTok{],}
\NormalTok{    xlabel}\OperatorTok{=}\StringTok{"antibiotic concentration"}\NormalTok{,}
\NormalTok{    ylabel}\OperatorTok{=}\StringTok{"optical density"}\NormalTok{,}
\NormalTok{    xscale}\OperatorTok{=}\NormalTok{log2}
\NormalTok{)}
\CommentTok{\# Plot data}
\FunctionTok{scatterlines!}\NormalTok{(ax, x, y)}

\NormalTok{fig}
\end{Highlighting}
\end{Shaded}

\includegraphics[width=7.29167in,height=6.25in]{iwasawa_bayesian_inference_files/figure-pdf/cell-23-output-1.png}

With the data simulated, let's perform inference on the data.

\begin{Shaded}
\begin{Highlighting}[]
\BuiltInTok{Random}\NormalTok{.}\FunctionTok{seed!}\NormalTok{(}\FloatTok{42}\NormalTok{)}
\CommentTok{\# Perform inference}
\NormalTok{model }\OperatorTok{=} \FunctionTok{logistic\_alt\_model}\NormalTok{(x, y)}

\CommentTok{\# Define number of steps}
\NormalTok{n\_burnin }\OperatorTok{=} \FloatTok{10\_000}
\NormalTok{n\_samples }\OperatorTok{=} \FloatTok{1\_000}

\NormalTok{chain }\OperatorTok{=}\NormalTok{ Turing.}\FunctionTok{sample}\NormalTok{(}
\NormalTok{    model, Turing.}\FunctionTok{NUTS}\NormalTok{(), Turing.}\FunctionTok{MCMCThreads}\NormalTok{(), n\_burnin }\OperatorTok{+}\NormalTok{ n\_samples, }\FloatTok{4}
\NormalTok{)}
\end{Highlighting}
\end{Shaded}

\begin{verbatim}
Sampling (4 threads)   0%|                              |  ETA: N/A
┌ Info: Found initial step size
└   ϵ = 0.4
┌ Info: Found initial step size
└   ϵ = 0.8
┌ Info: Found initial step size
└   ϵ = 1.6
┌ Info: Found initial step size
└   ϵ = 0.8
Sampling (4 threads)  25%|███████▌                      |  ETA: 0:00:21
Sampling (4 threads)  50%|███████████████               |  ETA: 0:00:07
Sampling (4 threads)  75%|██████████████████████▌       |  ETA: 0:00:02
Sampling (4 threads) 100%|██████████████████████████████| Time: 0:00:06
Sampling (4 threads) 100%|██████████████████████████████| Time: 0:00:06
\end{verbatim}

\begin{verbatim}
Chains MCMC chain (11000×17×4 Array{Float64, 3}):

Iterations        = 1001:1:12000
Number of chains  = 4
Samples per chain = 11000
Wall duration     = 5.85 seconds
Compute duration  = 22.85 seconds
parameters        = ic50, a, b, c, σ²
internals         = lp, n_steps, is_accept, acceptance_rate, log_density, hamiltonian_energy, hamiltonian_energy_error, max_hamiltonian_energy_error, tree_depth, numerical_error, step_size, nom_step_size

Summary Statistics
  parameters      mean       std      mcse     ess_bulk     ess_tail      rhat ⋯
      Symbol   Float64   Float64   Float64      Float64      Float64   Float64 ⋯

        ic50    1.3615    0.0783    0.0006   21996.9818   15834.2853    1.0001 ⋯
           a    0.9776    0.0641    0.0005   14338.2651   16755.4871    1.0000 ⋯
           b    7.5525    2.9801    0.0243   15179.8530   18757.0920    1.0000 ⋯
           c    0.1183    0.0456    0.0004   14321.6452   10118.2654    1.0002 ⋯
          σ²    0.0102    0.0066    0.0001   14265.0218   17759.7256    1.0002 ⋯
                                                                1 column omitted

Quantiles
  parameters      2.5%     25.0%     50.0%     75.0%     97.5%
      Symbol   Float64   Float64   Float64   Float64   Float64

        ic50    1.1980    1.3177    1.3639    1.4079    1.5098
           a    0.8545    0.9372    0.9761    1.0169    1.1094
           b    3.6070    5.6275    6.9973    8.7969   14.6326
           c    0.0281    0.0886    0.1180    0.1470    0.2112
          σ²    0.0037    0.0062    0.0085    0.0121    0.0270
\end{verbatim}

Let's look at the posterior distribution of the parameters.

\begin{Shaded}
\begin{Highlighting}[]
\CommentTok{\# Plot corner plot for chains}
\NormalTok{PairPlots.}\FunctionTok{pairplot}\NormalTok{(}
\NormalTok{    chain[n\_burnin}\OperatorTok{+}\FloatTok{1}\OperatorTok{:}\KeywordTok{end}\NormalTok{, }\OperatorTok{:}\NormalTok{, }\OperatorTok{:}\NormalTok{],}
\NormalTok{    PairPlots.}\FunctionTok{Truth}\NormalTok{(}
\NormalTok{        (;}
\NormalTok{        ic50}\OperatorTok{=}\NormalTok{ic50\_true,}
\NormalTok{        a}\OperatorTok{=}\NormalTok{a\_true,}
\NormalTok{        b}\OperatorTok{=}\NormalTok{b\_true,}
\NormalTok{        c}\OperatorTok{=}\NormalTok{c\_true,}
\NormalTok{        σ²}\OperatorTok{=}\NormalTok{σ²\_true}
\NormalTok{    )}
\NormalTok{    )}
\NormalTok{)}
\end{Highlighting}
\end{Shaded}

\includegraphics[width=19.20833in,height=19.60417in]{iwasawa_bayesian_inference_files/figure-pdf/cell-25-output-1.png}

The model is still able to recover the ground truth parameters. Let's
now plot the posterior predictive check.

\begin{Shaded}
\begin{Highlighting}[]
\BuiltInTok{Random}\NormalTok{.}\FunctionTok{seed!}\NormalTok{(}\FloatTok{42}\NormalTok{)}

\CommentTok{\# Initialize matrix to store samples}
\NormalTok{y\_samples }\OperatorTok{=} \FunctionTok{Array}\DataTypeTok{\{Float64\}}\NormalTok{(}\ConstantTok{undef}\NormalTok{, }\FunctionTok{length}\NormalTok{(x), n\_samples)}

\CommentTok{\# Loop through samples}
\ControlFlowTok{for}\NormalTok{ i }\KeywordTok{in} \FloatTok{1}\OperatorTok{:}\NormalTok{n\_samples}
    \CommentTok{\# Generate mean for sample}
\NormalTok{    y\_samples[}\OperatorTok{:}\NormalTok{, i] }\OperatorTok{=} \FunctionTok{logistic\_alt}\NormalTok{(}
\NormalTok{        x,}
\NormalTok{        chain[}\OperatorTok{:}\NormalTok{a][i],}
\NormalTok{        chain[}\OperatorTok{:}\NormalTok{b][i],}
\NormalTok{        chain[}\OperatorTok{:}\NormalTok{c][i],}
\NormalTok{        chain[}\OperatorTok{:}\NormalTok{ic50][i]}
\NormalTok{    )}
    \CommentTok{\# Add noise}
\NormalTok{    y\_samples[}\OperatorTok{:}\NormalTok{, i] }\OperatorTok{.+=} \FunctionTok{randn}\NormalTok{(}\FunctionTok{length}\NormalTok{(x)) }\OperatorTok{*} \FunctionTok{√}\NormalTok{(chain[}\OperatorTok{:}\NormalTok{σ²][i])}
\ControlFlowTok{end} \CommentTok{\# for}

\CommentTok{\# Initialize figure}
\NormalTok{fig }\OperatorTok{=} \FunctionTok{Figure}\NormalTok{(size}\OperatorTok{=}\NormalTok{(}\FloatTok{350}\NormalTok{, }\FloatTok{300}\NormalTok{))}
\CommentTok{\# Add axis}
\NormalTok{ax }\OperatorTok{=} \FunctionTok{Axis}\NormalTok{(}
\NormalTok{    fig[}\FloatTok{1}\NormalTok{, }\FloatTok{1}\NormalTok{],}
\NormalTok{    xlabel}\OperatorTok{=}\StringTok{"antibiotic concentration"}\NormalTok{,}
\NormalTok{    ylabel}\OperatorTok{=}\StringTok{"optical density"}\NormalTok{,}
\NormalTok{    title}\OperatorTok{=}\StringTok{"Posterior Predictive Check"}\NormalTok{,}
\NormalTok{    xscale}\OperatorTok{=}\NormalTok{log2}
\NormalTok{)}

\CommentTok{\# Plot samples}
\ControlFlowTok{for}\NormalTok{ i }\KeywordTok{in} \FloatTok{1}\OperatorTok{:}\NormalTok{n\_samples}
    \FunctionTok{lines!}\NormalTok{(ax, x, y\_samples[}\OperatorTok{:}\NormalTok{, i], color}\OperatorTok{=}\NormalTok{(}\OperatorTok{:}\NormalTok{gray, }\FloatTok{0.05}\NormalTok{))}
\ControlFlowTok{end} \CommentTok{\# for}

\CommentTok{\# Plot data}
\FunctionTok{scatterlines!}\NormalTok{(ax, x, y)}

\NormalTok{fig}
\end{Highlighting}
\end{Shaded}

\includegraphics[width=7.29167in,height=6.25in]{iwasawa_bayesian_inference_files/figure-pdf/cell-26-output-1.png}

Everything looks good. Let's again test it on the \textcite{iwasawa2022}
data.

\begin{Shaded}
\begin{Highlighting}[]
\BuiltInTok{Random}\NormalTok{.}\FunctionTok{seed!}\NormalTok{(}\FloatTok{42}\NormalTok{)}
\CommentTok{\# Define prior parameters}
\NormalTok{prior\_params }\OperatorTok{=}\NormalTok{ (}
\NormalTok{    a}\OperatorTok{=}\NormalTok{(}\FunctionTok{log}\NormalTok{(}\FloatTok{0.1}\NormalTok{), }\FloatTok{0.1}\NormalTok{),}
\NormalTok{    b}\OperatorTok{=}\NormalTok{(}\FloatTok{0}\NormalTok{, }\FloatTok{1}\NormalTok{),}
\NormalTok{    c}\OperatorTok{=}\NormalTok{(}\FloatTok{0}\NormalTok{, }\FloatTok{1}\NormalTok{),}
\NormalTok{    σ²}\OperatorTok{=}\NormalTok{(}\FloatTok{0}\NormalTok{, }\FloatTok{0.01}\NormalTok{)}
\NormalTok{)}
\CommentTok{\# Define data}
\NormalTok{data }\OperatorTok{=}\NormalTok{ df\_group[}\FloatTok{2}\NormalTok{]}
\CommentTok{\# Clean data}
\NormalTok{outliers\_idx }\OperatorTok{=} \FunctionTok{fit\_logistic\_and\_detect\_outliers}\NormalTok{(}
    \FunctionTok{log2}\NormalTok{.(data.concentration\_ugmL), data.OD, threshold}\OperatorTok{=}\FloatTok{2}
\NormalTok{)}
\NormalTok{data\_clean }\OperatorTok{=}\NormalTok{ data[.!outliers\_idx, }\OperatorTok{:}\NormalTok{]}

\CommentTok{\# Perform inference}
\NormalTok{model }\OperatorTok{=} \FunctionTok{logistic\_alt\_model}\NormalTok{(}
\NormalTok{    data\_clean.concentration\_ugmL,}
\NormalTok{    data\_clean.OD,}
\NormalTok{    prior\_params}
\NormalTok{)}

\CommentTok{\# Define number of steps}
\NormalTok{n\_burnin }\OperatorTok{=} \FloatTok{10\_000}
\NormalTok{n\_samples }\OperatorTok{=} \FloatTok{1\_000}

\NormalTok{chain }\OperatorTok{=}\NormalTok{ Turing.}\FunctionTok{sample}\NormalTok{(}
\NormalTok{    model, Turing.}\FunctionTok{NUTS}\NormalTok{(), Turing.}\FunctionTok{MCMCThreads}\NormalTok{(), n\_burnin }\OperatorTok{+}\NormalTok{ n\_samples, }\FloatTok{4}
\NormalTok{)}
\end{Highlighting}
\end{Shaded}

\begin{verbatim}
Sampling (4 threads)   0%|                              |  ETA: N/A
┌ Info: Found initial step size
└   ϵ = 0.00625
┌ Info: Found initial step size
└   ϵ = 0.0125
┌ Info: Found initial step size
└   ϵ = 0.00625
┌ Info: Found initial step size
└   ϵ = 0.00625
Sampling (4 threads)  25%|███████▌                      |  ETA: 0:00:21
Sampling (4 threads)  50%|███████████████               |  ETA: 0:00:07
Sampling (4 threads)  75%|██████████████████████▌       |  ETA: 0:00:02
Sampling (4 threads) 100%|██████████████████████████████| Time: 0:00:07
Sampling (4 threads) 100%|██████████████████████████████| Time: 0:00:07
\end{verbatim}

\begin{verbatim}
Chains MCMC chain (11000×17×4 Array{Float64, 3}):

Iterations        = 1001:1:12000
Number of chains  = 4
Samples per chain = 11000
Wall duration     = 6.52 seconds
Compute duration  = 24.81 seconds
parameters        = ic50, a, b, c, σ²
internals         = lp, n_steps, is_accept, acceptance_rate, log_density, hamiltonian_energy, hamiltonian_energy_error, max_hamiltonian_energy_error, tree_depth, numerical_error, step_size, nom_step_size

Summary Statistics
  parameters      mean       std      mcse     ess_bulk     ess_tail      rhat ⋯
      Symbol   Float64   Float64   Float64      Float64      Float64   Float64 ⋯

        ic50    8.5831    0.1769    0.0011   25568.1627   25963.0962    1.0003 ⋯
           a    0.1165    0.0038    0.0000   18041.1443   20901.4168    1.0002 ⋯
           b    9.6366    1.5316    0.0097   26259.3918   25025.7026    1.0000 ⋯
           c    0.0710    0.0035    0.0000   17294.7757   21118.9348    1.0002 ⋯
          σ²    0.0001    0.0000    0.0000   29855.3383   28397.5187    1.0001 ⋯
                                                                1 column omitted

Quantiles
  parameters      2.5%     25.0%     50.0%     75.0%     97.5%
      Symbol   Float64   Float64   Float64   Float64   Float64

        ic50    8.2417    8.4631    8.5828    8.6994    8.9339
           a    0.1091    0.1139    0.1164    0.1190    0.1241
           b    7.2053    8.5622    9.4488   10.4964   13.1592
           c    0.0640    0.0687    0.0711    0.0734    0.0777
          σ²    0.0001    0.0001    0.0001    0.0001    0.0002
\end{verbatim}

Let's look at the posterior distribution of the parameters.

\begin{Shaded}
\begin{Highlighting}[]
\CommentTok{\# Plot corner plot for chains}
\NormalTok{PairPlots.}\FunctionTok{pairplot}\NormalTok{(chain[n\_burnin}\OperatorTok{+}\FloatTok{1}\OperatorTok{:}\KeywordTok{end}\NormalTok{, }\OperatorTok{:}\NormalTok{, }\OperatorTok{:}\NormalTok{])}
\end{Highlighting}
\end{Shaded}

\includegraphics[width=19.5625in,height=20.02083in]{iwasawa_bayesian_inference_files/figure-pdf/cell-28-output-1.png}

Nothing looks obviously wrong. Let's plot the posterior predictive
check.

\begin{Shaded}
\begin{Highlighting}[]
\BuiltInTok{Random}\NormalTok{.}\FunctionTok{seed!}\NormalTok{(}\FloatTok{42}\NormalTok{)}

\CommentTok{\# Define unique concentrations}
\NormalTok{unique\_concentrations }\OperatorTok{=} \FunctionTok{sort}\NormalTok{(}\FunctionTok{unique}\NormalTok{(data\_clean.concentration\_ugmL))}

\CommentTok{\# Initialize matrix to store samples}
\NormalTok{y\_samples }\OperatorTok{=} \FunctionTok{Array}\DataTypeTok{\{Float64\}}\NormalTok{(}
    \ConstantTok{undef}\NormalTok{, }\FunctionTok{length}\NormalTok{(unique\_concentrations), n\_samples}
\NormalTok{)}

\CommentTok{\# Loop through samples}
\ControlFlowTok{for}\NormalTok{ i }\KeywordTok{in} \FloatTok{1}\OperatorTok{:}\NormalTok{n\_samples}
    \CommentTok{\# Generate mean for sample}
\NormalTok{    y\_samples[}\OperatorTok{:}\NormalTok{, i] }\OperatorTok{=} \FunctionTok{logistic\_alt}\NormalTok{(}
\NormalTok{        unique\_concentrations,}
\NormalTok{        chain[}\OperatorTok{:}\NormalTok{a][i],}
\NormalTok{        chain[}\OperatorTok{:}\NormalTok{b][i],}
\NormalTok{        chain[}\OperatorTok{:}\NormalTok{c][i],}
\NormalTok{        chain[}\OperatorTok{:}\NormalTok{ic50][i]}
\NormalTok{    )}
    \CommentTok{\# Add noise}
\NormalTok{    y\_samples[}\OperatorTok{:}\NormalTok{, i] }\OperatorTok{.+=} \FunctionTok{randn}\NormalTok{(}\FunctionTok{length}\NormalTok{(unique\_concentrations)) }\OperatorTok{*} \FunctionTok{√}\NormalTok{(chain[}\OperatorTok{:}\NormalTok{σ²][i])}
\ControlFlowTok{end} \CommentTok{\# for}

\CommentTok{\# Initialize figure}
\NormalTok{fig }\OperatorTok{=} \FunctionTok{Figure}\NormalTok{(size}\OperatorTok{=}\NormalTok{(}\FloatTok{350}\NormalTok{, }\FloatTok{300}\NormalTok{))}
\CommentTok{\# Add axis}
\NormalTok{ax }\OperatorTok{=} \FunctionTok{Axis}\NormalTok{(}
\NormalTok{    fig[}\FloatTok{1}\NormalTok{, }\FloatTok{1}\NormalTok{],}
\NormalTok{    xlabel}\OperatorTok{=}\StringTok{"antibiotic concentration"}\NormalTok{,}
\NormalTok{    ylabel}\OperatorTok{=}\StringTok{"optical density"}\NormalTok{,}
\NormalTok{    title}\OperatorTok{=}\StringTok{"Posterior Predictive Check"}\NormalTok{,}
\NormalTok{    xscale}\OperatorTok{=}\NormalTok{log2}
\NormalTok{)}

\CommentTok{\# Plot samples}
\ControlFlowTok{for}\NormalTok{ i }\KeywordTok{in} \FloatTok{1}\OperatorTok{:}\NormalTok{n\_samples}
    \FunctionTok{lines!}\NormalTok{(}
\NormalTok{        ax,}
\NormalTok{        unique\_concentrations,}
\NormalTok{        y\_samples[}\OperatorTok{:}\NormalTok{, i],}
\NormalTok{        color}\OperatorTok{=}\NormalTok{(ColorSchemes.Paired\_12[}\FloatTok{1}\NormalTok{], }\FloatTok{0.5}\NormalTok{)}
\NormalTok{    )}
\ControlFlowTok{end} \CommentTok{\# for}

\CommentTok{\# Plot data}
\FunctionTok{scatter!}\NormalTok{(ax, data\_clean.concentration\_ugmL, data\_clean.OD)}

\NormalTok{fig}
\end{Highlighting}
\end{Shaded}

\includegraphics[width=7.29167in,height=6.25in]{iwasawa_bayesian_inference_files/figure-pdf/cell-29-output-1.png}

With this parameterization, the model performs as well as the previous
one.

\subsection{Conclusion}\label{conclusion}

In this notebook, we have performed Bayesian inference on the data from
\textcite{iwasawa2022} using two different parameterizations of the
logistic function. Both parameterizations perform equally well, and we
are able to recover the ground truth parameters from simulated data.
Furthermore, implementing an outlier detection scheme improved the fit
significantly.



\end{document}
